\SourcePage{156}
\SGAsection{3}{拓扑斯态射的有限性条件}
\begin{SGAresult}{定义}{3.1}
设 $f:E'\to E$ 是\Term{topos-morphism}{拓扑斯态射}。若对 $E$ 的每个\Term{quasi-compact}{拟紧}（相应地，\Term{quasi-separated}{拟分离}）\Term{object}{对象} $X$，$f^*(X)$ 都\Term{quasi-compact}{拟紧}（相应地，\Term{quasi-separated}{拟分离}），则称 $f$ \Term{quasi-compact}{拟紧}（相应地，\Term{quasi-separated}{拟分离}）。若 $f$ 既\Term{quasi-compact}{拟紧}又\Term{quasi-separated}{拟分离}，则称 $f$ \Term{coherent}{凝聚}。
\end{SGAresult}
\textbf{3.1.1.} 注意，若 $f$ 具有上述性质之一，则与 $f$ \Term{isomorphism}{同构}的任何\Term{topos-morphism}{拓扑斯态射}也具有同一性质。两个\Term{quasi-compact}{拟紧}（相应地，\Term{quasi-separated}{拟分离}、\Term{coherent}{凝聚}）\Term{topos-morphism}{拓扑斯态射}的复合仍然\Term{quasi-compact}{拟紧}（相应地，\Term{quasi-separated}{拟分离}、\Term{coherent}{凝聚}），这一点也是显然的。

\begin{SGAresult}{命题}{3.2}
沿用3.1的记号，设 $(X_i)_{i\in I}$ 是 $E$ 中由\Term{quasi-compact}{拟紧}（相应地，\Term{coherent}{凝聚}）\Term{object}{对象}组成的\Term{generating-family}{生成族}。在括号中的相应情形下，假定 $E'$ 有由\Term{quasi-compact}{拟紧}\Term{object}{对象}组成的\Term{generating-family}{生成族}。$f$ \Term{quasi-compact}{拟紧}（相应地，\Term{coherent}{凝聚}）的充要条件是：对每个 $i\in I$，$f^*(X_i)$ 都\Term{quasi-compact}{拟紧}（相应地，\Term{coherent}{凝聚}）。
\end{SGAresult}
\begin{proof}
条件的必要性显然。为证第一种情形的充分性，注意对于 $E$ 的任一\Term{quasi-compact}{拟紧}\Term{object}{对象} $X$，有\Term{finite}{有限}\Term{epimorphic-family}{满态射族} $X_i\to X$，从而有\Term{finite}{有限}\Term{epimorphic-family}{满态射族} $f^*(X_i)\to f^*(X)$。依假设各 $f^*(X_i)$ \Term{quasi-compact}{拟紧}，所以 $f^*(X)$ \Term{quasi-compact}{拟紧}（1.3）。第二种情形中，还须说明：若 $X$ 是 $E$ 的\Term{quasi-separated}{拟分离}\Term{object}{对象}，则 $f^*(X)$ \Term{quasi-separated}{拟分离}。关于 $X$ 的假设可以表述为（1.17）存在\Term{epimorphic-family}{满态射族} $X_i\to X$，使各 $X_i\times_X X_j$ \Term{quasi-compact}{拟紧}。
\SourcePage{157}
于是有\Term{covering-family}{覆盖族} $f^*(X_i)\to f^*(X)$，且\Term{fiber-product}{纤维积}
\[
f^*(X_i)\mathbin{\times}_{f^*(X)}f^*(X_j)
\]
均\Term{quasi-compact}{拟紧}，因为由 a)，$f^*$ 把\Term{quasi-compact}{拟紧}\Term{object}{对象}变为\Term{quasi-compact}{拟紧}\Term{object}{对象}。各 $f^*(X_i)$ \Term{coherent}{凝聚}，因此再用1.17即得结论。
\end{proof}

\begin{SGAresult}{推论}{3.3}
设 $C,C'$ 是两个 $\mathcal U$-\Term{site}{景}，其中\Term{fiber-product}{纤维积}\Term{representable}{可表}，且每个\Term{covering-family}{覆盖族}都有\Term{finite-covering}{有限覆盖}子族；设 $g:C\to C'$ 是\Term{site-morphism}{景态射}（第IV讲4.9.1）。则由 $g$ 定义的\Term{topos-morphism}{拓扑斯态射} $f:C^{\sim}\to C'^{\sim}$ \Term{coherent}{凝聚}。
\end{SGAresult}
\begin{proof}
事实上，$\varepsilon_C(C)$（相应地，$\varepsilon_{C'}(C')$）是 $C$（相应地，$C'$）中满足3.2括号中相应条件的\Term{generating-family}{生成族}（2.1.1）。
\end{proof}

\begin{SGAresult}{命题}{3.4}
设 $E$ 是\Term{locally-coherent}{局部凝聚}\Term{topos}{拓扑斯}（2.3），$f:X\to Y$ 是 $E$ 的\Term{arrow}{箭头}，从而诱\Term{derivation}{导出}诱导\Term{topos}{拓扑斯}间的\Term{morphism}{态射}（第IV讲(5.5.2)）$E_{/X}\to E_{/Y}$。后者是\Term{quasi-compact}{拟紧}（相应地，\Term{quasi-separated}{拟分离}、\Term{coherent}{凝聚}）\Term{topos-morphism}{拓扑斯态射}的充要条件是：$f$ 是\Term{quasi-compact}{拟紧}（相应地，\Term{quasi-separated}{拟分离}、\Term{coherent}{凝聚}）\Term{morphism}{态射}。
\end{SGAresult}
\begin{proof}
“\Term{quasi-compact}{拟紧}”情形由定义即明（无需对 $E$ 加条件）；“\Term{quasi-separated}{拟分离}”情形就是2.8；“\Term{coherent}{凝聚}”情形则由前两种情形合起来得到。
\end{proof}

\begin{SGAresult}{命题}{3.5}
设 $f:E'\to E$ 是\Term{locally-coherent}{局部凝聚}\Term{topos}{拓扑斯}间的\Term{morphism}{态射}。设 $(X_i)_{i\in I}$ 是 $E$ 中由\Term{algebraic}{代数}\Term{coherent}{凝聚}\Term{object}{对象}组成的\Term{generating-family}{生成族}。下列条件\Term{equivalence}{等价}：

(i) 对 $E$ 的每条\Term{quasi-compact}{拟紧}\Term{arrow}{箭头} $u$，$f^*(u)$ 都\Term{quasi-compact}{拟紧}。

(i bis) 对每条\Term{arrow}{箭头} $u:X_i\to X_j$，$f^*(u)$ 都\Term{quasi-compact}{拟紧}。

(ii) 对 $E'$ 的每个\Term{coherent}{凝聚}\Term{object}{对象} $Y'$、$E$ 的每个\Term{algebraic}{代数}\Term{coherent}{凝聚}\Term{object}{对象} $Y$，及每条\Term{arrow}{箭头} $v:Y'\to f^*(Y)$，相应的\Term{topos-morphism}{拓扑斯态射}
\[
f_v:E'_{/Y'}\longrightarrow E_{/Y},
\]
即由 $Y$ \Term{derivation}{导出}的\Term{localization}{局部化}\Term{morphism}{态射} $E'_{/Y'}\to E'_{/f^*(Y)}$（第IV讲(5.5.2)）与由 $f$ 诱导的\Term{morphism}{态射} $E'_{/f^*(Y)}\to E_{/Y}$（第IV讲(5.10.1)）的复合，是\Term{coherent}{凝聚}的。

(ii bis) 与(ii)相同，但只考虑一组资料 $v_\alpha:Y'_\alpha\to f^*(Y_\alpha)$，使 $Y'_\alpha$ 都是\Term{algebraic}{代数}的，并\Term{covering}{覆盖} $E'$ 的\Term{final-object}{终对象}。

(ii ter)（给定 $E'$ 中由\Term{coherent}{凝聚}\Term{object}{对象}组成的\Term{generating-family}{生成族} $(X'_{i'})_{i'\in I'}$，并对每个 $i'\in I'$ 给定一个 $i\in I$ 及\Term{arrow}{箭头} $v_{i'}:X'_{i'}\to f(X_i)$。）对每个 $i'\in I'$，由 $v_{i'}$ 定义的\Term{topos-morphism}{拓扑斯态射} $E'_{/X'_{i'}}\to E_{/X_i}$ \Term{coherent}{凝聚}。
\end{SGAresult}
\begin{proof}
每条\Term{arrow}{箭头} $X_i\to X_j$ 都\Term{quasi-compact}{拟紧}，故显然(i)$\Rightarrow$(i bis)。反之，假定(i bis)成立，来证明(i)。设 $u:X\to Y$ 是 $E$ 的\Term{quasi-compact}{拟紧}\Term{arrow}{箭头}。利用\Term{epimorphic-family}{满态射族} $X_i\to Y$，关于 $u$ 的假设\Term{representable}{可表}述为各 $X\times_Y X_i$ \Term{quasi-compact}{拟紧}（1.16），即对每个 $i$，存在\Term{finite}{有限}\Term{epimorphic-family}{满态射族} $X_j\to X\times_Y X_i$（1.3）。
\SourcePage{158}
于是有相应的\Term{epimorphic-family}{满态射族} $f^*(X_i)\to f^*(Y)$，且对每个 $i$，有\Term{finite}{有限}\Term{epimorphic-family}{满态射族}
\[
f^*(X_j)\longrightarrow f^*(X\times_Y X_i)
\simeq f^*(X)\mathbin{\times}_{f^*(Y)}f^*(X_i).
\]
由(i bis)，这些\Term{arrow}{箭头}与到 $f^*(X_i)$ 的\Term{projection}{投影} $\operatorname{pr}_2$ 的复合是\Term{quasi-compact}{拟紧}\Term{morphism}{态射} $f^*(X_j)\to f^*(X_i)$。因此（1.11(i)），对每个 $i$，$\operatorname{pr}_2:f^*(X)\times_{f^*(Y)}f^*(X_i)$ \Term{quasi-compact}{拟紧}，故（1.10(ii)）$f^*(u):f^*(X)\to f^*(Y)$ \Term{quasi-compact}{拟紧}。

还须证明(i)$\Rightarrow$(ii)和(ii bis)$\Rightarrow$(i)；(ii)$\Rightarrow$(ii bis)是显然的，而(ii ter)是(ii bis)的特例。(i)$\Rightarrow$(ii)直接成立：事实上，$E_{/Y}$ 的\Term{object}{对象} $X$ \Term{quasi-compact}{拟紧}（相应地，\Term{quasi-separated}{拟分离}），当且仅当结构\Term{morphism}{态射} $X\to Y$（相应地，对角\Term{morphism}{态射} $X\to X\times_Y X$）\Term{quasi-compact}{拟紧}（2.7）；对于 $f_Y^*(X')$，也有同样的\Term{quasi-compact}{拟紧}性和\Term{quasi-separated}{拟分离}性判据。最后证明(ii bis)$\Rightarrow$(i)。设 $u:X\to Y$ 是 $E$ 的\Term{quasi-compact}{拟紧}\Term{arrow}{箭头}，须证 $f^*(u)$ \Term{quasi-compact}{拟紧}。由于 $Y'_\alpha$ \Term{covering}{覆盖} $E'$ 的\Term{final-object}{终对象}，只需证明各 $f^*(u)|_{Y'_\alpha}$ \Term{quasi-compact}{拟紧}（1.10(ii)）。这些\Term{morphism}{态射}正是 $f_{v_\alpha}^*(u|_{Y_\alpha})$，所以归结于证明下面的结论。
\end{proof}

\begin{SGAresult}{推论}{3.6}
设 $f:E'\to E$ 是\Term{locally-coherent}{局部凝聚}\Term{topos}{拓扑斯}间的\Term{coherent}{凝聚}\Term{morphism}{态射}。则对 $E$ 的每条\Term{quasi-compact}{拟紧}\Term{arrow}{箭头} $u$，$f^*(u)$ 都\Term{quasi-compact}{拟紧}。
\end{SGAresult}
\begin{proof}
由已经证明的3.4中的蕴涵(i bis)$\Rightarrow$(i)，只须证明：若 $u:X\to Y$ 是 $E$ 的\Term{arrow}{箭头}，$X,Y$ \Term{coherent}{凝聚}，则 $f^*(u)$ 是\Term{quasi-compact}{拟紧}\Term{morphism}{态射}。这立即由关于 $f$ 的假设得到，因为该假设蕴涵 $f^*(X)$ 与 $f^*(Y)$ 也\Term{coherent}{凝聚}。
\end{proof}

\begin{SGAresult}{定义}{3.7}
设 $f:E'\to E$ 是\Term{locally-coherent}{局部凝聚}\Term{topos}{拓扑斯}间的\Term{morphism}{态射}。若 $f$ 满足3.5中的\Term{equivalence}{等价}条件，则称 $f$ \Term{locally-coherent}{局部凝聚}。
\end{SGAresult}
注意，这里的“\Term{local-scheme}{局部}”意指在上面\Term{local-scheme}{局部}。

\textbf{3.7.1.} 这一概念仍只依赖于\Term{topos-morphism}{拓扑斯态射} $f$ 的\Term{isomorphism}{同构}类，并且显然对\Term{topos-morphism}{拓扑斯态射}的复合\Term{stable}{稳定}。此外，由3.5，若 $f$ \Term{coherent}{凝聚}，则它\Term{locally-coherent}{局部凝聚}；若 $E',E$ 都\Term{coherent}{凝聚}，则逆命题成立（由3.5(ii bis)）。由判据3.5(ii bis)立即得到下述判据。

\begin{SGAresult}{推论}{3.7}
设 $f:E'\to E$ 是\Term{locally-coherent}{局部凝聚}\Term{topos}{拓扑斯}间的\Term{morphism}{态射}，$(Y_i)_{i\in I}$、$(Y'_i)_{i\in I}$ 分别是 $E$、$E'$ 的两族\Term{object}{对象}；并对每个 $i\in I$，设有\Term{morphism}{态射} $v_i:Y'_i\to f^*(Y_i)$，从而有\Term{topos-morphism}{拓扑斯态射}
\[
f_{v_i}:E'_{/Y'_i}\longrightarrow E_{/Y_i}.
\]
\SourcePage{159}
假定 $Y'_i$ \Term{covering}{覆盖} $E'$ 的\Term{final-object}{终对象}。则 $f$ \Term{locally-coherent}{局部凝聚}，当且仅当对每个 $i\in I$，$f_{v_i}$ 都\Term{locally-coherent}{局部凝聚}。
\end{SGAresult}
特别地，将其应用于诱导\Term{topos}{拓扑斯}的恒等\Term{morphism}{态射}，得到：
\begin{SGAresult}{推论}{3.8}
设 $E$ 是\Term{locally-coherent}{局部凝聚}\Term{topos}{拓扑斯}，$v:X\to Y$ 是 $E$ 的\Term{arrow}{箭头}。则由 $v$ 定义的诱导\Term{topos-morphism}{拓扑斯态射} $E_{/X}\to E_{/Y}$ \Term{locally-coherent}{局部凝聚}。
\end{SGAresult}

\begin{SGAresult}{注记}{3.9}
似乎实践中遇到的所有\Term{algebraic}{代数}\Term{topos-morphism}{拓扑斯态射}都是\Term{locally-coherent}{局部凝聚}的（所以“\Term{locally-coherent}{局部凝聚}”一词大概不会得到非常广泛的使用）。这也就是说，实践中遇到的\Term{coherent}{凝聚}\Term{topos}{拓扑斯}间的\Term{morphism}{态射}都是\Term{coherent}{凝聚}的。不过，可以构造\Term{coherent}{凝聚}\Term{topos}{拓扑斯}间不\Term{coherent}{凝聚}的\Term{morphism}{态射}；更具体地，可以构造从\Term{point-topos}{点拓扑斯}（第IV讲2.2）到某个\Term{noetherian}{诺特}\Term{topos}{拓扑斯} $E$ 的不\Term{coherent}{凝聚}\Term{morphism}{态射}，见2.17 i)。
\end{SGAresult}

\begin{SGAresult}{例}{3.10}
（\Term{scheme}{概形}所关联的\Term{topos}{拓扑斯}。）回到1.22中配备SGA IV 6.3的\Term{topology}{拓扑} $T_i$ 的\Term{category}{范畴} $(\mathrm{Sch})$ 的例子。对每个\Term{scheme}{概形} $X$，设 $X_{T_i}$ 为诱导\Term{site}{景} $(\mathrm{Sch})_{/X}$ 所定义的 $\mathcal V$-\Term{topos}{拓扑斯}。\Term{scheme}{概形}\Term{morphism}{态射} $f:X\to Y$ 定义\Term{topos-morphism}{拓扑斯态射} $f_{T_i}:X_{T_i}\to Y_{T_i}$；由3.4和1.22，后者\Term{quasi-compact}{拟紧}（相应地，\Term{quasi-separated}{拟分离}、\Term{coherent}{凝聚}），当且仅当\Term{scheme}{概形}\Term{morphism}{态射} $f$ 在EGA IV 1的通常意义下\Term{quasi-compact}{拟紧}（相应地，\Term{quasi-separated}{拟分离}、\Term{coherent}{凝聚}）。无论如何，由3.8，$f_{T_i}$ \Term{locally-coherent}{局部凝聚}。

也可以像1.22.1中那样，为\Term{scheme}{概形} $X$ 关联 $\mathcal U$-\Term{topos}{拓扑斯} $\operatorname{Top}(X)$ 与 $\operatorname{Top}(X_{\text{ét}})$；每个\Term{scheme}{概形}\Term{morphism}{态射} $f:X\to Y$ 因而关联相应\Term{topos}{拓扑斯}间的\Term{morphism}{态射} $\operatorname{Top}(f)$ 与 $\operatorname{Top}(f_{\text{ét}})$。以 $T(f)$ 表示这两个\Term{topos-morphism}{拓扑斯态射}中的任一个。利用3.2立即验证，此\Term{topos-morphism}{拓扑斯态射}总是\Term{locally-coherent}{局部凝聚}的，并且它\Term{quasi-compact}{拟紧}（相应地，\Term{quasi-separated}{拟分离}、\Term{coherent}{凝聚}），当且仅当\Term{scheme}{概形}\Term{morphism}{态射} $f$ 在通常意义下\Term{quasi-compact}{拟紧}（相应地，\Term{quasi-separated}{拟分离}、\Term{coherent}{凝聚}）。

这些观察再次表明，本节引入的术语与\Term{scheme}{概形}论中通用的术语相容，因而即使最不情愿的读者，也应接受它们。
\end{SGAresult}

\begin{SGAresult}{习题}{3.11}
（\Term{coherent}{凝聚}\Term{topos}{拓扑斯}与\Term{pretopos}{预拓扑斯}。）

a) 称满足下列条件的\Term{category}{范畴} $C$ 为 $\mathcal U$-\Term{pretopos}{预拓扑斯}（prétopos，或简称\Term{pretopos}{预拓扑斯}）：

1) $C$ 中的\Term{finite}{有限}\Term{inverse-limit}{射影极限}\Term{representable}{可表}。

2) $C$ 中的\Term{finite}{有限}和\Term{representable}{可表}，且它们是不交的、\Term{universal}{普遍}的。
\SourcePage{160}
3) $C$ 中的\Term{equivalence-relation}{等价关系}都是\Term{effective}{有效的}，且 $C$ 中每个\Term{epimorphism}{满态射}都是\Term{universal}{普遍}\Term{effective}{有效的}。

4) $C$ \Term{equivalence}{等价}于某个属于 $\mathcal U$ 的\Term{category}{范畴}。

证明：若 $E$ 是\Term{coherent}{凝聚} $\mathcal U$-\Term{topos}{拓扑斯}，则由 $E$ 的\Term{coherent}{凝聚}\Term{object}{对象}组成的\Term{full-subcategory}{满子范畴} $E_{\mathrm{coh}}$ 是 $\mathcal U$-\Term{pretopos}{预拓扑斯}，且包含\Term{functor}{函子} $E_{\mathrm{coh}}\to E$ \Term{left-exact}{左正合}，与\Term{finite}{有限}和及按\Term{equivalence-relation}{等价关系}取商交换。（用1.5.3、1.15、1.17.1。）证明：$E$ 在 $C$ 上诱导的\Term{topology}{拓扑}是“\Term{precanonical-topology}{预典范}”（précanonique）\Term{topology}{拓扑}，即其\Term{covering-family}{覆盖族} $X_i\to X$ 是那些含有关于 $C$ 的\Term{canonical-topology}{典范拓扑}为\Term{covering}{覆盖}的\Term{finite}{有限}子族的族。因此，知道\Term{pretopos}{预拓扑斯} $C=E_{\mathrm{coh}}$，便可将 $E$ 重建为\Term{topos}{拓扑斯} $C^{\sim}$，准确到\Term{equivalence}{等价}，其中 $C$ 配备\Term{precanonical-topology}{预典范}\Term{topology}{拓扑}。

b) 设 $C$ 是 $\mathcal U$-\Term{pretopos}{预拓扑斯}。给 $C$ 配备\Term{precanonical-topology}{预典范}\Term{topology}{拓扑}。证明：$C$ 的每个\Term{morphism}{态射} $f$ 都\Term{decomposable}{可分解}为 $f''f'$，其中 $f'$ 是\Term{epimorphism}{满态射}，$f''$ 是\Term{monomorphism}{单态射}。证明：\Term{topos}{拓扑斯} $E=C^{\sim}$ \Term{coherent}{凝聚}，且\Term{canonical}{典范}\Term{functor}{函子} $\varepsilon:C\to E$ 诱导 $C$ 与 $E$ 的子\Term{category}{范畴} $E_{\mathrm{coh}}$ 的\Term{equivalence}{等价}。（先证明 $\varepsilon$ 是\Term{fully-faithful}{全忠实}\Term{left-exact}{左正合}\Term{functor}{函子}，与\Term{finite}{有限}和及按\Term{equivalence-relation}{等价关系}取商交换；再证明，$E$ 中一个 $C$ 的\Term{object}{对象}的\Term{coherent}{凝聚}\Term{subobject}{子对象}属于 $C$ 的\Term{essential-image}{本质像}。）所以，只要知道关联的 $\mathcal U$-\Term{topos}{拓扑斯} $E=C^{\sim}$，便可重建 $\mathcal U$-\Term{pretopos}{预拓扑斯} $C$，准确到\Term{category}{范畴}\Term{equivalence}{等价}。

c) 设 $C,C'$ 是两个 $\mathcal U$-\Term{pretopos}{预拓扑斯}。证明：\Term{functor}{函子} $\varphi:C'\to C$ 是从 $C$ 到 $C'$ 的\Term{site-morphism}{景态射}（关于\Term{precanonical-topology}{预典范}\Term{topology}{拓扑}）的充要条件是：$\varphi$ \Term{left-exact}{左正合}，与\Term{finite}{有限}和及按\Term{equivalence-relation}{等价关系}取商交换。

d) 沿用 c) 的记号，假定 $C,C'$ 如 a) 中那样分别关联于\Term{coherent}{凝聚} $\mathcal U$-\Term{topos}{拓扑斯} $E,E'$。证明\Term{canonical}{典范}\Term{functor}{函子}（第IV讲4.9.3）
\[
\operatorname{Mor}_{\mathrm{site}}(C,C')\longrightarrow
\mathscr{H}\!om_{\mathrm{top}}(E,E')
\]
诱导第一个\Term{category}{范畴}（已在 c) 中明确写出）与第二个\Term{category}{范畴}中由\Term{coherent}{凝聚}\Term{topos-morphism}{拓扑斯态射} $E\to E'$ 组成的\Term{full-subcategory}{满子范畴}之间的\Term{equivalence}{等价}；一个\Term{quasi-inverse}{拟逆}\Term{functor}{函子}由下述方法得到：对每个\Term{coherent}{凝聚}\Term{topos-morphism}{拓扑斯态射} $f:E\to E'$，关联由 $F^*$ 诱导的\Term{functor}{函子} $E'_{\mathrm{coh}}=C'\to E_{\mathrm{coh}}=C$。

e) 设 $C$ 是 $\mathcal U$-\Term{pretopos}{预拓扑斯}，$E=C^{\sim}$。直接用 $C$ 的\Term{exactness}{正合性}性质表述1.25中的\Term{equivalence}{等价}条件（参见1.28 b)）。证明：$C$ 的\Term{object}{对象} $X$ 在 $E$ 中是\Term{noetherian}{诺特}的，当且仅当 $C$ 中 $X$ 的每个递增\Term{subobject}{子对象}列都\Term{stable}{稳定}；故 $E$ \Term{noetherian}{诺特}，当且仅当 $C$ 的每个\Term{object}{对象}满足上述条件。
\end{SGAresult}

\begin{SGAresult}{习题}{3.12}
若存在\Term{finite}{有限}\Term{category}{范畴} $C$ 使\Term{topos}{拓扑斯} $E$ \Term{equivalence}{等价}于 $\widehat C$，则称 $E$ 为\Term{finite-topos}{有限拓扑斯}。
\SourcePage{161}
a) $E$ 为\Term{finite-topos}{有限拓扑斯}的充要条件是：$E$ 有足够多\Term{essential-point}{本质点}（第IV讲7.6 b)），且 $E$ 的\Term{essential-point}{本质点}\Term{category}{范畴} $\operatorname{Pointess}(E)$ \Term{equivalence}{等价}于\Term{finite}{有限}\Term{category}{范畴}。（用第IV讲7.6 d)、h)。）

b) 假定 $E$ \Term{finite}{有限}。证明 $E$ 的每个点都是\Term{essential-point}{本质点}。证明2-\Term{functor}{函子} $E\mapsto\operatorname{Point}(E)$ 与 $C\mapsto\widehat C$ 在\Term{finite-topos}{有限拓扑斯} $E$ 的2-\Term{category}{范畴}及\Term{equivalence}{等价}于\Term{finite}{有限}\Term{category}{范畴}的\Term{karoubian-category}{Karoubi范畴}（catégorie karoubienne）$C$（第IV讲7.5 a)）的2-\Term{category}{范畴}之间建立2-\Term{equivalence}{等价}，并且\Term{finite-topos}{有限拓扑斯}间的每个\Term{morphism}{态射}都是本质的（第IV讲7.6 a)）。（用 a) 和第IV讲7.6 h)。）

c) 假定 $E=\widehat C$，其中 $C$ \Term{equivalence}{等价}于\Term{finite}{有限}\Term{category}{范畴}，设 $F$ 是\Term{coherent}{凝聚}\Term{topos}{拓扑斯}。回忆（第IV讲(4.6.3.1)），\Term{topos-morphism}{拓扑斯态射} $f:E\to F$ 对应于\Term{functor}{函子} $u:C\to\operatorname{Point}(F)$。证明：$f$ \Term{coherent}{凝聚}，当且仅当 $f$ 把 $E$ 的\Term{coherent}{凝聚}点（3.1）变为 $F$ 的\Term{coherent}{凝聚}点，亦即当且仅当对每个 $X\in\operatorname{Ob}C$，$u(X)$ 都是 $F$ 的\Term{coherent}{凝聚}点。（用3.2和2.17 g)。）由此推出，从\Term{finite-topos}{有限拓扑斯}到\Term{finite-topos}{有限拓扑斯}的每个\Term{morphism}{态射} $f:E\to F$ 都\Term{coherent}{凝聚}。

d) 设 $X$ 是\Term{sober-space}{清醒拓扑空间}（第IV讲4.2.1）。\Term{topos}{拓扑斯} $\operatorname{Top}(X)$（第IV讲2.1）\Term{finite}{有限}的充要条件是：$X$ 为\Term{finite}{有限}集。（用第IV讲7.1.6。）
\end{SGAresult}

\SGAsection{4}{由粘合得到的拓扑斯中的有限性条件}
本节在本研讨班后续部分不再使用。

\textbf{4.1.} 设 $E$ 是\Term{topos}{拓扑斯}，$U$ 是 $E$ 的\Term{open-object}{开对象}（即 $E$ 的\Term{final-object}{终对象}的\Term{subobject}{子对象}），考虑 $E$ 中相应的\Term{open-subtopos}{开子拓扑斯}和\Term{closed-subtopos}{闭子拓扑斯}（第IV讲9）
\[
E'=E_{/U},\qquad E''=E_{\complement U}.
\]
以
\[
j:E'\longrightarrow E,\qquad i:E''\longrightarrow E
\]
表示\Term{canonical-topology}{典范拓扑}斯\Term{morphism}{态射}。回忆（3.4），$j$ 是\Term{quasi-compact}{拟紧}（相应地，\Term{coherent}{凝聚}）\Term{topos-morphism}{拓扑斯态射}的充要条件是：$U$ 到 $E$ 的\Term{final-object}{终对象}的包含
\[
j:U\longrightarrow e
\]
（也记作 $j$）是 $E$ 中的\Term{quasi-compact}{拟紧}（相应地，\Term{coherent}{凝聚}）\Term{morphism}{态射}；另一方面，$j:E'\to E$ 总是\Term{quasi-separated}{拟分离}的，因为 $j:U\to e$ \Term{quasi-separated}{拟分离}（1.8.1）。我们拟给出\Term{topos}{拓扑斯} $E''$ \Term{coherent}{凝聚}、\Term{topos-morphism}{拓扑斯态射} $i:E''\to E$ \Term{coherent}{凝聚}的判据。首先指出：

\begin{SGAresult}{命题}{4.2}
沿用4.1的记号，包含\Term{morphism}{态射} $i:E''\to E$ \Term{quasi-compact}{拟紧}，即对 $E$ 的每个\Term{quasi-compact}{拟紧}\Term{object}{对象} $X$，$E''$ 的\Term{object}{对象} $i^*(X)$ \Term{quasi-compact}{拟紧}；此外，若 $X$ \Term{prenoetherian}{预诺特}（1.30），则 $i^*(X)$ \Term{prenoetherian}{预诺特}。
\end{SGAresult}
这由定义1.1和下述引理得到。

\SourcePage{162}
\begin{SGAresult}{引理}{4.2.1}
映射 $Y\mapsto i^*(Y)$ 在 $X$ 的包含\Term{subobject}{子对象} $X_U=X\times U$ 的\Term{subobject}{子对象}集，与 $i^*(X)$ 的\Term{subobject}{子对象}集之间，诱导\Term{ordered-set}{有序集}\Term{isomorphism}{同构}。
\end{SGAresult}
\begin{proof}
注意，$i_*$ \Term{conservative}{保守}（因为它\Term{fully-faithful}{全忠实}）且\Term{left-exact}{左正合}（特别地，与\Term{fiber-product}{纤维积}交换），故 $E''$ 中的\Term{morphism}{态射} $u:Y\to Z$ 是\Term{monomorphism}{单态射}，当且仅当 $i_*(u):i_*(Y)\to i_*(Z)$ 是\Term{monomorphism}{单态射}。由此立即得到，利用 $i_*$ 把 $E'$ 与 $E$ 中由满足第IV讲9.7 1)中\Term{equivalence}{等价}条件的\Term{object}{对象} $Z$ 组成的\Term{full-subcategory}{满子范畴}等同后，$E''$ 的\Term{object}{对象} $Z$ 在 $E'$ 中的\Term{subobject}{子对象}，等同于 $Z$ 在 $E$ 中那些属于 $E''$ 的\Term{subobject}{子对象} $Y$；也就是说，那些包含 $Z_U\xrightarrow{\sim}U$ 的\Term{subobject}{子对象}。对于形如 $i_*i^*(X)=X_{\complement U}$ 的 $Z$，由第IV讲(9.5.1)给出的\Term{pushout}{并合和}
\[
i_*i^*(X)=X_{\complement U}=X\amalg_{X_U}U,
\]
给定其在 $E$ 中的一个\Term{subobject}{子对象} $Y$，\Term{equivalence}{等价}于给定 $X$ 的\Term{subobject}{子对象} $Y_1$ 和 $U$ 的\Term{subobject}{子对象} $Y_2$，使这两个\Term{subobject}{子对象}在 $X_U$ 中的\Term{inverse-image}{逆像}相同。条件 $Y$ 包含 $U$ 意味着 $Y_2=U$，而对 $Y_1$ 剩下的条件是 $Y_1$ 包含 $X_U$。这就建立了 $i^*(X)$ 在 $E''$ 中的\Term{subobject}{子对象}集，与 $X$ 的包含 $X_U$ 的\Term{subobject}{子对象}集之间的\Term{bijection}{双射}。显然这是\Term{ordered-set}{有序集}\Term{isomorphism}{同构}，且逆映射正是3.14.1中所说的映射。
\end{proof}

\begin{SGAresult}{推论}{4.3}
a) 设 $X$ 是 $E$ 的\Term{object}{对象}。$i^*(X)$ 和 $j^*(X)$ \Term{quasi-compact}{拟紧}，是 $X$ \Term{quasi-compact}{拟紧}的充分条件；若 $j:U\to e$ \Term{quasi-compact}{拟紧}，则这也是必要条件。

b) 设 $Y$ 是 $E''$ 的\Term{object}{对象}。$i_*(Y)$ \Term{quasi-compact}{拟紧}，是 $Y$ \Term{quasi-compact}{拟紧}的充分条件；若 $U$ \Term{quasi-compact}{拟紧}，则这也是必要条件。
\end{SGAresult}
\begin{proof}
a) 充分性由定义1.1及\Term{functor}{函子}对 $(i^*,j^*)$ \Term{conservative}{保守}这一事实（第IV讲9.11 3)）得到。必要性由 $i,j$ \Term{quasi-compact}{拟紧}得到（由3.14及 $j$ \Term{quasi-compact}{拟紧}的假设）。

b) 因为 $Y\simeq i^*i_*(Y)$，充分性由4.2得到。必要性由 a) 中的充分性得到，注意 $i^*i_*(Y)\simeq Y$ 且 $j^*i_*(Y)\simeq U$。
\end{proof}

\begin{SGAresult}{推论}{4.4}
设 $Y$ 是 $E''$ 的\Term{object}{对象}。若 $U$ \Term{quasi-compact}{拟紧}，或 $E$ 有由\Term{quasi-compact}{拟紧}\Term{object}{对象}组成的\Term{generating-family}{生成族}，则 $i_*(Y)$ \Term{quasi-separated}{拟分离}（相应地，\Term{coherent}{凝聚}）是 $Y$ \Term{quasi-separated}{拟分离}（相应地，\Term{coherent}{凝聚}）的充分条件。若 $U$ \Term{quasi-separated}{拟分离}（相应地，\Term{coherent}{凝聚}），且 $j:U\to e$ 为\Term{quasi-compact}{拟紧}\Term{morphism}{态射}，则 $i_*(Y)$ \Term{quasi-separated}{拟分离}（相应地，\Term{coherent}{凝聚}）是 $Y$ \Term{quasi-separated}{拟分离}（相应地，\Term{coherent}{凝聚}）的必要条件。
\end{SGAresult}
\begin{proof}
由4.3 b)，括号中的相应情形由不带括号的情形得到。假定 $i_*(Y)$ \Term{quasi-separated}{拟分离}，来证明在所作的两个假设之一成立时，$Y$ 也\Term{quasi-separated}{拟分离}。
\SourcePage{163}
须证明：对 $Y$ 上方的两个\Term{quasi-compact}{拟紧}\Term{object}{对象} $Y',Y''$，\Term{fiber-product}{纤维积} $Y'\times_Y Y''$ \Term{quasi-compact}{拟紧}。当假定 $Y$ \Term{quasi-compact}{拟紧}时，注意只须证明 $i_*(Y'\times_Y Y'')$ \Term{quasi-compact}{拟紧}（4.3 b)）；这是由于 $i_*(Y'),i_*(Y'')$ \Term{quasi-compact}{拟紧}（4.3 b)，这里用了 $U$ 的\Term{quasi-compact}{拟紧}性），$i_*$ 与\Term{fiber-product}{纤维积}交换，以及 $i_*(Y)$ \Term{quasi-separated}{拟分离}的假设。当假定 $E$ 有由\Term{quasi-compact}{拟紧}\Term{object}{对象}组成的\Term{generating-family}{生成族}时，$i_*(Y')$ 是其\Term{quasi-compact}{拟紧}\Term{subobject}{子对象} $X'_\alpha$ 的\Term{filtered-direct-limit}{滤过归纳极限}，故 $Y'\simeq i^*i_*(Y')$ 是\Term{subobject}{子对象} $i^*(X'_\alpha)$ 的\Term{filtered-direct-limit}{滤过归纳极限}；又因 $Y'$ \Term{quasi-compact}{拟紧}，它等于某个 $i^*(X'_\alpha)$。同样，$Y''$ 形如 $i^*(X''_\beta)$，其中 $X''_\beta$ 是 $i_*(Y'')$ 的\Term{quasi-compact}{拟紧}\Term{subobject}{子对象}。于是
\[
Y'\times_Y Y''=i^*\bigl(X'_\alpha\mathbin{\times}_{i_*Y}X''_\beta\bigr).
\]
由假设 $i_*(Y)$ \Term{quasi-separated}{拟分离}，$X'_\alpha\times_{i_*(Y)}X''_\beta$ \Term{quasi-compact}{拟紧}，再由4.2，$Y'\times_Y Y''$ 也\Term{quasi-compact}{拟紧}。

反之，假定 $U$ \Term{quasi-separated}{拟分离}且 $j:U\to e$ \Term{quasi-compact}{拟紧}，来证明若 $Y$ \Term{quasi-separated}{拟分离}，则 $i_*(Y)$ 也\Term{quasi-separated}{拟分离}；即对于 $i_*(Y)$ 上方的两个\Term{quasi-compact}{拟紧}\Term{object}{对象} $X',X''$，\Term{fiber-product}{纤维积} $X'\times_{i_*(Y)}X''$ \Term{quasi-compact}{拟紧}。应用4.3 a)，只须证明它在 $i^*$ 下的像\Term{quasi-compact}{拟紧}（这由4.2及 $Y$ \Term{quasi-separated}{拟分离}的假设得到），并证明它在 $j^*$ 下的像\Term{quasi-compact}{拟紧}。而后者是 $j^*(X')\times j^*(X'')$，它确实\Term{quasi-compact}{拟紧}，因为 $j^*(X'),j^*(X'')$ \Term{quasi-compact}{拟紧}（$j$ \Term{quasi-compact}{拟紧}），且 $U$ \Term{quasi-separated}{拟分离}。
\end{proof}

\begin{SGAresult}{推论}{4.5}
假定 $E$ 有由\Term{quasi-compact}{拟紧}（相应地，\Term{prenoetherian}{预诺特}）\Term{object}{对象}组成的\Term{generating-subcategory}{生成子范畴}，则 $E''$ 也如此。
\end{SGAresult}
这由4.2及下述引理得到。
\begin{SGAresult}{引理}{4.5.1}
若 $(X_\alpha)_\alpha$ 是 $E$ 中的\Term{generating-family}{生成族}，则 $(i^*(X_\alpha))_\alpha$ 是 $E''$ 中的\Term{generating-family}{生成族}。
\end{SGAresult}
\begin{proof}
这实际上意味着 $E''$ 上的\Term{functor}{函子}族
\[
Y\longmapsto\operatorname{Hom}(i^*(X_\alpha),Y)
\]
\Term{conservative}{保守}。由于 $\operatorname{Hom}(i^*(X_\alpha),Y)\simeq\operatorname{Hom}(X_\alpha,i_*(Y))$，只须用 $i_*$ \Term{conservative}{保守}这一事实。
\end{proof}

\begin{SGAresult}{命题}{4.6}
记号同4.1。

a) 若 $E$ 有由\Term{coherent}{凝聚}\Term{object}{对象}组成的\Term{generating-family}{生成族}，则\Term{closed-subtopos}{闭子拓扑斯} $E''$ 也如此，且包含\Term{morphism}{态射} $i:E''\to E$ \Term{coherent}{凝聚}。

b) 假定 $E$ 中的\Term{morphism}{态射} $j:U\to e$ \Term{quasi-compact}{拟紧}。若 $E$ \Term{locally-coherent}{局部凝聚}（相应地，\Term{coherent}{凝聚}、\Term{quasi-separated}{拟分离}、\Term{locally-noetherian}{局部诺特}、\Term{noetherian}{诺特}、\Term{locally-perfect}{局部完美}、\Term{perfect}{完美}（parfait）），则 $E''$ 也如此，且包含\Term{morphism}{态射} $i:E''\to E$ \Term{coherent}{凝聚}。
\end{SGAresult}
\SourcePage{164}
\begin{proof}
a) 第一个断言将由第二个断言及4.5.1得到。无论如何，4.5蕴涵 $E''$ 有由\Term{quasi-compact}{拟紧}\Term{object}{对象}组成的\Term{generating-family}{生成族}，所以可用3.2；它表明只须验证，对于 $E$ 的每个\Term{coherent}{凝聚}\Term{object}{对象} $X$，$E''$ 的\Term{object}{对象} $i^*(X)$ \Term{coherent}{凝聚}。利用\Term{open-object}{开对象}的补\Term{topos}{拓扑斯}的构造与\Term{localization}{局部化}相容（第IV讲9.13 b)），可归结于 $X$ 是 $E$ 的\Term{final-object}{终对象}的情形；所以须证明 $E''$ 的\Term{final-object}{终对象}\Term{coherent}{凝聚}。这时可用4.4，只须证明 $i_*(e_{E''})=e_E$ \Term{coherent}{凝聚}，而这一点成立。

b) 在 b) 所作的每个假设下，$E$ 都有由\Term{coherent}{凝聚}\Term{object}{对象}组成的\Term{generating-family}{生成族}，所以由 a) 已知 $i:E''\to E$ \Term{coherent}{凝聚}。再由4.5.1，$E''$ 中由\Term{coherent}{凝聚}\Term{object}{对象}组成的\Term{full-subcategory}{满子范畴} $E''_{\mathrm{coh}}$ 是\Term{generating-subcategory}{生成子范畴}。假定 $E$ \Term{coherent}{凝聚}，来证明 $E''$ 也\Term{coherent}{凝聚}；即（2.4.5）须由 $E$ 的子\Term{category}{范畴} $E_{\mathrm{coh}}$ 对\Term{finite}{有限}\Term{inverse-limit}{射影极限}\Term{stable}{稳定}，证明 $E''$ 的子\Term{category}{范畴} $E''_{\mathrm{coh}}$ 也如此。这立即由4.4括号中的相应判据及 $i_*$ \Term{left-exact}{左正合}得到；该判据可以应用，因为现在 $U$ \Term{coherent}{凝聚}，这是由于 $e$ \Term{coherent}{凝聚}且 $j:U\to e$ \Term{quasi-compact}{拟紧}。通过\Term{localization}{局部化}，再用第IV讲9.13 b)，便得到若 $E$ \Term{locally-coherent}{局部凝聚}，则 $E''$ 也\Term{locally-coherent}{局部凝聚}。若 $E$ \Term{quasi-separated}{拟分离}，则 $E''$ 也\Term{quasi-separated}{拟分离}：事实上，由4.4，$E''$ 的\Term{final-object}{终对象}\Term{quasi-separated}{拟分离}，因为 $E$ 的\Term{final-object}{终对象}如此；还须证明 $E''$ 的两个\Term{quasi-separated}{拟分离}\Term{object}{对象}的积\Term{quasi-separated}{拟分离}，已知在 $E''$ 中有这一性质。这仍由判据4.4及 $i_*$ 与积交换得到；注意 $e$ 的\Term{subobject}{子对象} $U$ \Term{quasi-separated}{拟分离}，因为 $e$ 如此，所以4.4适用。

一个\Term{topos}{拓扑斯}\Term{locally-noetherian}{局部诺特}（相应地，\Term{noetherian}{诺特}），当且仅当它\Term{locally-coherent}{局部凝聚}（相应地，\Term{coherent}{凝聚}）且有由\Term{prenoetherian}{预诺特}\Term{object}{对象}组成的\Term{generating-family}{生成族}（2.10(iii bis)）；所以由上述结论及4.5，若 $E$ \Term{locally-noetherian}{局部诺特}（相应地，\Term{noetherian}{诺特}），则 $E''$ 也如此。现在假定 $E$ \Term{perfect}{完美}，来证明 $E''$ \Term{perfect}{完美}。由于已经知道它\Term{coherent}{凝聚}，只须验证它满足判据1.25(i)，即 $E''$ 的每个\Term{object}{对象} $Y$ 都是\Term{coherent}{凝聚}\Term{object}{对象}的\Term{filtered-direct-limit}{滤过归纳极限}，已知相应断言在 $E$ 中成立。而 $i_*(Y)$ 是\Term{coherent}{凝聚}\Term{object}{对象} $X_\alpha$ 的\Term{filtered-direct-limit}{滤过归纳极限}，故 $Y\simeq i^*i_*(Y)$ 是\Term{coherent}{凝聚}\Term{object}{对象} $i^*(X_\alpha)$ 的\Term{filtered-direct-limit}{滤过归纳极限}。通过\Term{localization}{局部化}（第IV讲9.13 b)），由此得到若 $E$ \Term{locally-perfect}{局部完美}，则 $E''$ 也\Term{locally-perfect}{局部完美}。4.6证毕。
\end{proof}

\begin{SGAresult}{注记}{4.6.1}
事实上，在4.6 b)中，$j:U\to e$ \Term{quasi-compact}{拟紧}的假设是不必要的，惟“$E$ \Term{quasi-separated}{拟分离}”情形可能例外。由上述证明，只须说明 $E$ \Term{coherent}{凝聚}蕴涵 $E''$ \Term{coherent}{凝聚}。$U$ 是其\Term{quasi-compact}{拟紧}\Term{subobject}{子对象} $U_\alpha$ 的\Term{filtered-direct-limit}{滤过归纳极限}；在第7节中将验证，这时 $E''$ 等同于 $E_{/U_i}$ 的补\Term{closed-subtopos}{闭子拓扑斯} $E''_i$ 的\Term{inverse-limit}{射影极限}\Term{topos}{拓扑斯}（在第7节的意义下）。由4.6 b)，这些都是\Term{coherent}{凝聚}\Term{topos}{拓扑斯}，且过渡\Term{morphism}{态射}\Term{coherent}{凝聚}。
\SourcePage{165}
因此 $E''$ 是\Term{coherent}{凝聚}\Term{topos}{拓扑斯}（第7节）。
\end{SGAresult}

\begin{SGAresult}{推论}{4.7}
假定 $E$ 中由\Term{coherent}{凝聚}\Term{object}{对象}组成的子\Term{category}{范畴} $E_{\mathrm{coh}}$ 是生成的，设 $X$ 是 $E$ 的\Term{object}{对象}。

a) $i^*(X)$ 和 $j^*(X)$ \Term{quasi-separated}{拟分离}，是 $X$ \Term{quasi-separated}{拟分离}的必要条件；若 $j:U\to e$ \Term{quasi-compact}{拟紧}，则这也是充分条件。

b) 假定 $j:U\to e$ \Term{quasi-compact}{拟紧}。则 $X$ \Term{coherent}{凝聚}，当且仅当 $i^*(X)$ 和 $j^*(X)$ \Term{coherent}{凝聚}。

c) 假定 $E$ \Term{coherent}{凝聚}且 $j:U\to e$ \Term{quasi-compact}{拟紧}。则 $E$ \Term{perfect}{完美}（2.9.1），当且仅当 $E'$ 和 $E''$ \Term{perfect}{完美}。
\end{SGAresult}
\begin{proof}
a) 在4.1中已指出 $j:E'\to E$ \Term{quasi-separated}{拟分离}；另一方面，由4.6 a)，$i$ 也\Term{quasi-separated}{拟分离}，故必要性成立。为证充分性，设 $X',X''$ 是 $X$ 上方的\Term{quasi-compact}{拟紧}\Term{object}{对象}，须证 $X'\times_X X''$ \Term{quasi-compact}{拟紧}；只须证明它在 $i^*$、$j^*$ 下的像\Term{quasi-compact}{拟紧}（4.3 a)）。由4.2及 $j:U\to e$ \Term{quasi-compact}{拟紧}的假设，$i,j$ 都\Term{quasi-compact}{拟紧}，故利用 $i^*,j^*$ 与\Term{fiber-product}{纤维积}交换，以及 $i^*(X),j^*(X)$ \Term{quasi-separated}{拟分离}的假设，即得结论。

b) 将 a) 与4.3 a) 合起来即得。

c) 必要性已经证明（4.6 b)）。为证充分性，注意 $E$ \Term{coherent}{凝聚}，故 $E',E''$ 都\Term{coherent}{凝聚}（4.6 b)）；对于所考虑的每个\Term{topos}{拓扑斯}，它\Term{perfect}{完美}意味着其\Term{coherent}{凝聚}\Term{object}{对象}\Term{category}{范畴}对\Term{finite}{有限}\Term{direct-limit}{归纳极限}\Term{stable}{稳定}。于是由 b) 得到结论。
\end{proof}

\begin{SGAresult}{习题}{4.8}
a) 解决下述问题（撰稿人惭愧地承认，他不知道答案）：沿用4.1的记号，若 $E$ \Term{quasi-separated}{拟分离}（相应地，\Term{algebraic}{代数}），$E''$ 是否也如此？

b) 假定 $E$ \Term{equivalence}{等价}于形如 $\widehat C$ 的\Term{topos}{拓扑斯}，其中\Term{category}{范畴} $C$ \Term{equivalence}{等价}于某个属于 $\mathcal U$ 的\Term{category}{范畴}。利用2.17 c)及第IV讲9.24，直接证明：若 $E$ \Term{locally-coherent}{局部凝聚}（相应地，\Term{coherent}{凝聚}、\Term{algebraic}{代数}、\Term{quasi-separated}{拟分离}、\Term{locally-noetherian}{局部诺特}、\Term{noetherian}{诺特}），则 $E''$ 也如此。
\end{SGAresult}

\textbf{4.9.} 沿用4.1的记号。回忆（第IV讲9.16），给定由\Term{topos}{拓扑斯} $E$ 及 $E$ 的\Term{open-object}{开对象} $U$ 组成的情形 $(E,U)$，本质上\Term{equivalence}{等价}于给定三元组 $(E',E'',f)$，其中 $E',E''$ 是\Term{topos}{拓扑斯}，且
\[
f:E'\longrightarrow E''\tag{4.9.1}
\]
是“\Term{gluing-functor}{粘合函子}”，即\Term{left-exact}{左正合}且\Term{accessible}{可达}（accessible）的\Term{functor}{函子}。从4.1中的 $(E,U)$ 出发，相应的\Term{gluing-functor}{粘合函子}为
\[
f=i^*j_* .\tag{4.9.2}
\]
\SourcePage{166}
我们拟用资料 $E',E'',f$ 表述下述条件：$E$ 是\Term{coherent}{凝聚}（相应地，\Term{perfect}{完美}）\Term{topos}{拓扑斯}，且 $U$ 是\Term{quasi-compact}{拟紧}\Term{object}{对象}；\Term{equivalence}{等价}地，$E$ 与 $E'$ \Term{coherent}{凝聚}。已经知道这蕴涵 $E''$ 也\Term{coherent}{凝聚}（4.6 b)），因此问题归结为：给定两个\Term{coherent}{凝聚}\Term{topos}{拓扑斯} $E',E''$ 及\Term{gluing-functor}{粘合函子} $f$（4.9.1），$f$ 满足什么条件时，\Term{gluing}{粘合}所得的\Term{topos}{拓扑斯} $E$ \Term{coherent}{凝聚}（相应地，\Term{perfect}{完美}）？此外，已知（4.7 c)）$E$ \Term{perfect}{完美}，当且仅当 $E$ \Term{coherent}{凝聚}且 $E',E''$ \Term{perfect}{完美}，所以括号中的相应问题可归结为不带括号的问题。不过指出，若已假定 $E'$ \Term{perfect}{完美}，后一个问题的解答（4.10）会更漂亮。

\textbf{4.9.3.} 首先注意，若 $E$ \Term{coherent}{凝聚}，则 $E$ 中由\Term{coherent}{凝聚}\Term{object}{对象} $X=(X',X'',u:X''\to f(X'))$ 组成的\Term{full-subcategory}{满子范畴} $E_{\mathrm{coh}}$，可重建为 $E$ 中由满足 $X'=j^*(X)$、$X''=i^*(X)$ \Term{coherent}{凝聚}的\Term{object}{对象} $X$ 组成的子\Term{category}{范畴} $E_0$（4.7 b)）。所以 $E$ \Term{coherent}{凝聚}的一个必要条件是上述子\Term{category}{范畴} $E_0$ 生成。这也是充分条件，因为由4.3 a)，$E_0$ 由\Term{quasi-compact}{拟紧}\Term{object}{对象}组成；另一方面，$E_0$ 显然对\Term{finite}{有限}\Term{inverse-limit}{射影极限}\Term{stable}{稳定}，再用3.4.5即得结论。

现在回忆（第IV讲9.18），给定\Term{gluing-functor}{粘合函子}(4.9.1)，\Term{equivalence}{等价}于给定 $E''$ 上取值于 $\operatorname{Pro}(E')^{\circ}$ 的\Term{sheaf}{层}：
\[
G\in\operatorname{Faisc}(E'',\operatorname{Pro}(E')^{\circ}),\quad
G:E''^{\circ}\longrightarrow\operatorname{Pro}(E')^{\circ};\tag{4.9.4}
\]
\Term{equivalence}{等价}地，给定与\Term{direct-limit}{归纳极限}交换的\Term{functor}{函子} $g=G^{\circ}$，
\[
g:E''\longrightarrow\operatorname{Pro}(E').\tag{4.9.5}
\]
若 $C$ 是 $E''$ 的\Term{generating-subcategory}{生成子范畴}，给它配备\Term{induced-topology}{诱导拓扑}，则已知（第II讲6.10）给定 $G$ 也\Term{equivalence}{等价}于给定 $C$ 上取值于 $\operatorname{Pro}(E')^{\circ}$ 的\Term{sheaf}{层}（准确到唯一\Term{isomorphism}{同构}）
\[
G_C\in\operatorname{Faisc}(C,\operatorname{Pro}(E')^{\circ}),\quad
G_C:C^{\circ}\longrightarrow\operatorname{Pro}(E')^{\circ};\tag{4.9.6}
\]
\Term{equivalence}{等价}地，给定满足已知的\Term{left-exact}{左正合}性条件（第II讲6.2.1）的\Term{functor}{函子} $g_C=G_C^{\circ}$，
\[
g_C:C\longrightarrow\operatorname{Pro}(E').\tag{4.9.7}
\]
当然，$g_C$ 就是 $g$（4.9.5）在 $C$ 上的\Term{restriction}{限制}。对我们最有意义的是取 $C=E''_{\mathrm{coh}}$ 的情形，因而得到\Term{object}{对象}
\[
G_0\in\operatorname{Faisc}(E''_{\mathrm{coh}},\operatorname{Pro}(E')^{\circ}),\quad
G_0:E''_{\mathrm{coh}}{}^{\circ}\longrightarrow\operatorname{Pro}(E')^{\circ},\tag{4.9.8}
\]
\[
g_0=G_0^{\circ}:E''_{\mathrm{coh}}\longrightarrow\operatorname{Pro}(E'),\tag{4.9.9}
\]
其中每一个仍\Term{equivalence}{等价}于给定 $f$。

\SourcePage{167}
\textbf{4.9.10.} 假定 $E''$ 的满\Term{generating-subcategory}{生成子范畴} $C$ 由\Term{quasi-compact}{拟紧}\Term{object}{对象}组成，并且在 $E''$ 中对\Term{finite}{有限}和、按\Term{equivalence-relation}{等价关系}取商及\Term{fiber-product}{纤维积}\Term{stable}{稳定}；例如 $C=E''_{\mathrm{coh}}$ 时如此（1.15及1.17.1）。那么，若 $P$ 是\Term{finite}{有限}\Term{inverse-limit}{射影极限}\Term{representable}{可表}的\Term{category}{范畴}（例如 $P=\operatorname{Pro}(E')^{\circ}$），便立即看出：\Term{functor}{函子} $G_C:G^{\circ}\to P$ 是取值于 $P$ 的\Term{sheaf}{层}，当且仅当\Term{functor}{函子} $g_C=G_C^{\circ}:C\to P^{\circ}$ 与\Term{finite}{有限}和及按\Term{equivalence-relation}{等价关系}取商交换。这特别明确了哪些是(4.9.6)中的\Term{sheaf}{层}，因而也表述了\Term{gluing-functor}{粘合函子} $f:E'\to E''$。

作了这些回顾以后，可以给出4.9中问题的解答。

\begin{SGAresult}{命题}{4.10}
设 $E',E''$ 是两个\Term{coherent}{凝聚}\Term{topos}{拓扑斯}，$f:E'\to E''$ 是\Term{gluing-functor}{粘合函子}（第IV讲9.10），即\Term{left-exact}{左正合}且\Term{accessible}{可达}的\Term{functor}{函子}。设 $E$ 是由其\Term{gluing}{粘合}得到的\Term{topos}{拓扑斯}（第IV讲9.16），以 $E'_{\mathrm{coh}}$（相应地，$E'_{\mathrm{PF}}$）表示 $E'$ 中由\Term{coherent}{凝聚}\Term{object}{对象}（相应地，由满足其所代表的\Term{covariant}{协变}\Term{functor}{函子} $\operatorname{Hom}(X,-)$ 与小\Term{filtered-direct-limit}{滤过归纳极限}交换的\Term{object}{对象} $X$）组成的\Term{strictly-full-subcategory}{严格满子范畴}。考虑下列条件：

(i) $E$ \Term{coherent}{凝聚}。

(ii) 对 $E'$ 的每个\Term{object}{对象} $X'$，存在一族以 $X'$ 为目标、以 $E'$ 的\Term{coherent}{凝聚}\Term{object}{对象}为源的\Term{morphism}{态射}
\[
v_\alpha:Y'_\alpha\longrightarrow X',
\]
使得 $E''$ 中的族
\[
f(v_\alpha):f(Y'_\alpha)\longrightarrow f(X')
\]
为\Term{covering-family}{覆盖族}。

(ii bis) $f$ 与（小）\Term{filtered-direct-limit}{滤过归纳极限}交换，且(ii)对每个 $X'\in\operatorname{Ob}E'_{\mathrm{PF}}$ 成立。

(ii ter) 由\Term{gluing-functor}{粘合函子} $f$ 定义的\Term{canonical}{典范}\Term{functor}{函子}（第IV讲9.20.1）
\[
\pi:\operatorname{Point}(E'')\longrightarrow\operatorname{Pro}(E')\tag{4.10.1}
\]
通过 $\operatorname{Pro}(E'_{\mathrm{coh}})$ \Term{resolution}{分解}（准确到\Term{isomorphism}{同构}）；这里使用包含\Term{functor}{函子} $E'_{\mathrm{coh}}\to E'$ 所给出的\Term{fully-faithful}{全忠实}\Term{functor}{函子} $\operatorname{Pro}(E'_{\mathrm{coh}})\to\operatorname{Pro}(E')$：
\[
\pi_0:\operatorname{Point}(E'')\longrightarrow\operatorname{Pro}(E'_{\mathrm{coh}}).\tag{4.10.2}
\]

(iii) \Term{functor}{函子} $f$ 与（小）\Term{filtered-direct-limit}{滤过归纳极限}交换。

(iii bis) \Term{functor}{函子} $g_0$（4.9.9）通过\Term{functor}{函子}
\[
E''_{\mathrm{coh}}\longrightarrow\operatorname{Pro}(E'_{\mathrm{PF}})\tag{4.10.3}
\]
\Term{resolution}{分解}（准确到\Term{isomorphism}{同构}），这里使用包含\Term{functor}{函子} $E'_{\mathrm{PF}}\to E'$ 所\Term{derivation}{导出}的\Term{fully-faithful}{全忠实}\Term{functor}{函子} $\operatorname{Pro}(E'_{\mathrm{PF}})\to\operatorname{Pro}(E')$。
\SourcePage{168}
(iii ter) \Term{canonical}{典范}\Term{functor}{函子} $\pi$（4.10.1）\Term{resolution}{分解}为（准确到\Term{isomorphism}{同构}）
\[
\pi_1:\operatorname{Point}(E'')\longrightarrow\operatorname{Pro}(E'_{\mathrm{PF}}).\tag{4.10.4}
\]

(iv) \Term{functor}{函子} $g_0$（4.9.9）通过\Term{functor}{函子}
\[
E''_{\mathrm{coh}}\longrightarrow\operatorname{Pro}(E'_{\mathrm{coh}})\tag{4.10.5}
\]
\Term{resolution}{分解}（准确到\Term{isomorphism}{同构}）。

则有如下蕴涵图：
\[
\begin{gathered}
\text{(iv)}\Rightarrow\text{(i)}\Leftrightarrow\text{(ii)}\Leftrightarrow\text{(ii bis)}\Leftrightarrow\text{(ii ter)}\\
\Rightarrow\text{(iii)}\Leftrightarrow\text{(iii bis)}\Leftrightarrow\text{(iii ter)}.
\end{gathered}
\tag{4.10.6}
\]
\end{SGAresult}
立即指出下述推论。
\begin{SGAresult}{推论}{4.11}
a) 假定 $E'$ \Term{perfect}{完美}。则4.10中考虑的所有条件\Term{equivalence}{等价}；特别地，$E$ \Term{coherent}{凝聚}，当且仅当 $f$ 与\Term{filtered-direct-limit}{滤过归纳极限}交换，亦即当且仅当\Term{functor}{函子} $g_0$（4.9.9）通过 $\operatorname{Pro}(E'_{\mathrm{coh}})$ \Term{resolution}{分解}（准确到\Term{isomorphism}{同构}）。

b) $E$ \Term{perfect}{完美}的充要条件是：$E',E''$ \Term{perfect}{完美}，且 $f$（相应地，$g_0$）满足 a) 中的条件。
\end{SGAresult}
\begin{proof}
事实上，$E'$ \Term{perfect}{完美}意味着 $E'_{\mathrm{coh}}=E'_{\mathrm{PF}}$，所以(iii bis)蕴涵(iv)，由此得到 a)。如4.9的预备说明所示，断言 b) 随之得到。
\end{proof}

\textbf{4.12. 4.10的证明。}

a) 明确写出4.9.3中得到的与(i)\Term{equivalence}{等价}的条件，即 $E$ 的\Term{full-subcategory}{满子范畴} $E_0$ 生成；这就是说，对 $E$ 的每个\Term{object}{对象} $X=(X',X'',u:X''\to f(X'))$，所有形如
\[
v=(v',v''):Y=(Y',Y'',v:Y''\longrightarrow f(y'))\longrightarrow X
\]
而 $Y\in\operatorname{Ob}E_0$（即 $Y',Y''$ \Term{coherent}{凝聚}）的\Term{morphism}{态射}所成的族是\Term{epimorphic-family}{满态射族}。\Term{functor}{函子}对 $(i^*,j^*)$ \Term{conservative}{保守}，所以这也意味着相应\Term{morphism}{态射}族 $Y'\to X'$ 在 $E'$ 中为\Term{epimorphic-family}{满态射族}，且相应\Term{morphism}{态射}族 $Y''\to X''$ 在 $E''$ 中为\Term{epimorphic-family}{满态射族}。对第一个族，这一点显然；只需取 $U$ 上方的 $Y\in\operatorname{Ob}E_0$，即满足 $Y''=\varnothing_{E''}$ 的 $Y$，便得到 $E'$ 中所有以 $X'$ 为目标、以\Term{coherent}{凝聚}\Term{object}{对象}为源的\Term{arrow}{箭头}所成的族。它确为\Term{epimorphic-family}{满态射族}，因为 $E'$ \Term{coherent}{凝聚}，故子\Term{category}{范畴} $E'_{\mathrm{coh}}$ 生成。所以只须表述：对任意给定的 $E$ 的\Term{object}{对象} $X$，族 $Y''\to X''$ 是\Term{epimorphic-family}{满态射族}。

b) 假定(ii)对这里考虑的\Term{object}{对象} $X'$ 成立。于是可找到\Term{epimorphic-family}{满态射族} $Y''_\beta\to X''$，其中每个\Term{morphism}{态射}与 $u:X''\to f(X')$ 的复合都\Term{lifting}{提升}为一个\Term{morphism}{态射} $Y''_\beta\to f(Y'_\alpha)$，这里取适当的 $\alpha=\varphi(\beta)$。
\SourcePage{169}
由于 $E''_{\mathrm{coh}}$ 是 $E''$ 的\Term{generating-subcategory}{生成子范畴}，甚至可以假定各 $Y''_\beta$ \Term{coherent}{凝聚}。于是\Term{commutative-diagram}{交换图}
\[
\begin{tikzcd}
Y''_\beta\arrow[r]\arrow[d]&f(Y'_{\varphi(\beta)})\arrow[d]\\
X''\arrow[r]&f(X')
\end{tikzcd}
\]
给出所求的 $E$ 中的\Term{morphism}{态射}族，并给出\Term{epimorphic-family}{满态射族} $Y''_\beta\to X''$。所以(ii)$\Rightarrow$(i)。

反之，假定(i)，来证明(ii)。把 a) 中明确写出的条件应用于\Term{object}{对象}
\[
X=j_*(X')=(X',f(X'),\operatorname{id}=f(X')\xrightarrow{\sim}f(X')).
\]
由于 $Y''\to X''\simeq f(X')$ 组成\Term{epimorphic-family}{满态射族}，族 $f(v'):f(Y')\to f(X')$ 更是\Term{epimorphic-family}{满态射族}；又依假设 $Y'$ \Term{coherent}{凝聚}，故得结论。所以
\[
\text{(i)}\Longleftrightarrow\text{(ii)}.
\]

c) 现在注意，条件(iv)也意味着：对于 $F''$ 的每个\Term{object}{对象} $X''$，$E'$ 上的\Term{pro-representable}{亲可表}（pro-représentable）\Term{functor}{函子}
\[
g(X''):X'\longmapsto\operatorname{Hom}(X'',f(X'))
\]
可由各\Term{set-component}{分量}属于 $E'_{\mathrm{coh}}$ 的\Term{pro-object}{亲对象}（pro-objet）\Term{pro-represent}{亲表示}；\Term{equivalence}{等价}地，在\Term{category}{范畴} $E'_{/g(X'')}$ 中存在\Term{cofinal}{余尾}\Term{object}{对象}集 $(X'_\alpha,u_\alpha)$，其中 $u_\alpha:X''\to f(X'_\alpha)$，$X''_\alpha\in\operatorname{Ob}E'_{\mathrm{coh}}$。这里 $E'_{/g(X'')}$ 的\Term{object}{对象}是一个 $E'$ 的\Term{object}{对象} $X'$ 与一个元素 $u\in g(X'')(X')=\operatorname{Hom}(X'',f(X'))$ 所成的对，即 $u:X''\to f(X')$。由于\Term{category}{范畴} $E'_{/g(X'')}$ \Term{filtered}{滤过}，\Term{cofinal}{余尾}性只意味着：对每个\Term{object}{对象} $(X',u)$，$u:X''\to f(X')$，都存在从某个 $(X'_i,u_i)$ 到它的\Term{morphism}{态射}；即存在\Term{morphism}{态射} $v'_\alpha:X'_\alpha\to X'$，使三角图
\[
\begin{tikzcd}
&f(x'_\alpha)\arrow[d,"f'(v'_\alpha)"]\\
X''\arrow[ur,"u_\alpha"]\arrow[r]&f(X')
\end{tikzcd}\tag{*}
\]
交换。所以条件(iv)\Term{equivalence}{等价}于：

(iv bis) 对每条\Term{arrow}{箭头} $u:X''\to f(X')$，其中 $X'\in\operatorname{Ob}E'$、$X''\in\operatorname{Ob}E'_{\mathrm{coh}}$，存在 $E'$ 中的\Term{arrow}{箭头} $v:Y'\to X'$，使 $u$ 通过 $f(v'):f(X')\to f(X')$ \Term{resolution}{分解}。

\SourcePage{170}
显然，条件(iv bis)蕴涵(ii)，因为对固定的 $X'$，以\Term{coherent}{凝聚}\Term{object}{对象}为源的所有 $u:X''\to f(X')$ 组成\Term{epimorphic-family}{满态射族}；从而，相应图(*)中所有 $f(v'_\alpha)$ 所成的族更是\Term{epimorphic-family}{满态射族}。所以
\[
\text{(iv)}\Rightarrow\text{(i)}.
\]

d) 明确写出条件(ii ter)。设 $p$ 是 $E''$ 的点；按定义，$\pi(p)$ 是\Term{pro-represent}{亲表示}\Term{functor}{函子} $h:X'\mapsto f(X')_p$ 的 $E'$ 的\Term{pro-object}{亲对象}；亦即由\Term{canonical}{典范}\Term{functor}{函子} $E'_{/h}\to E'$ 定义的\Term{pro-object}{亲对象}，其中 $E'_{/h}$ 表示由 $X'\in\operatorname{Ob}E'$ 和 $u\in h(X')$（即 $u\in f(X')_p$）所组成的对 $(X',u)$ 的\Term{category}{范畴}。这个\Term{pro-object}{亲对象}\Term{isomorphism}{同构}于来自 $\operatorname{Pro}(E'_{\mathrm{coh}})$ 的\Term{pro-object}{亲对象}，意味着 $E'_{/h}$ 中由满足 $X'\in\operatorname{Ob}E'_{\mathrm{coh}}$ 的 $(X',u)$ 组成的子\Term{category}{范畴}是\Term{cofinal}{余尾}的；即对于 $E'_{/h}$ 的每个\Term{object}{对象} $(X',u)$，$u\in f(X')_p$，存在\Term{dominant}{支配}它的\Term{object}{对象} $(X'_0,u_0)$，且 $X'_0$ \Term{coherent}{凝聚}；也就是存在\Term{morphism}{态射} $X'_0\to X'$，其中 $X'_0$ \Term{coherent}{凝聚}，使
\[
u\in\operatorname{Im}\bigl(f(X'_0)\longrightarrow f(X')\bigr).
\]
这对每个点 $p$ 成立，意味着对 $E'$ 的每个\Term{object}{对象} $X'$，族 $f(X'_0)\to f(X')$ 所给出的 $f(X'_0)_P\to f(X')_P$ \Term{surjectivity}{满射}。如将在附录（第10节）中看到的，\Term{coherent}{凝聚}\Term{topos}{拓扑斯} $E''$ 有\Term{enough-points}{足够多点}，故所得条件也意味着对每个 $X'\in\operatorname{Ob}E$，族 $f(X'_0)\to f(X')$ 是\Term{epimorphic-family}{满态射族}。这正是条件(ii)，所以
\[
\text{(ii)}\Longleftrightarrow\text{(ii ter)}.
\]

e) $f$ 与\Term{filtered-direct-limit}{滤过归纳极限}交换的条件（条件(iii)），由于 $E''$ 的\Term{fiber-functor}{纤维函子}族\Term{conservative}{保守}，\Term{equivalence}{等价}于复合\Term{functor}{函子} $h:X'\mapsto f(X')_P$ 与\Term{filtered-direct-limit}{滤过归纳极限}交换。来证明这意味着\Term{pro-represent}{亲表示}该\Term{functor}{函子}的\Term{pro-object}{亲对象} $\pi(p)$ 属于 $\operatorname{Pro}(E'_{\mathrm{PF}})$ 的\Term{essential-image}{本质像}；这便证明了蕴涵
\[
\text{(iii)}\Leftrightarrow\text{(iii ter)}\Leftarrow\text{(ii ter)}\Rightarrow\text{(iii bis)}.
\]
所以我们须证明下述引理。

\begin{SGAresult}{引理}{4.12.1}
设 $E'$ 是\Term{topos}{拓扑斯}，其子\Term{category}{范畴} $E'_{\mathrm{coh}}$ 是生成的，且 $X'\in\operatorname{Ob}\operatorname{Pro}(E')$。$X'$ 属于 $\operatorname{Pro}(E'_{\mathrm{PF}})$ 的\Term{essential-image}{本质像}，当且仅当它所\Term{pro-represent}{亲表示}的\Term{functor}{函子} $h:E'\to(\mathrm{Ens})$ 与\Term{filtered-direct-limit}{滤过归纳极限}交换。
\end{SGAresult}
\begin{proof}
与\Term{filtered-direct-limit}{滤过归纳极限}交换的\Term{functor}{函子}的\Term{filtered-direct-limit}{滤过归纳极限}仍具有同一性质，故充分性由 $E'_{\mathrm{PF}}$ 的定义得到。为证必要性，须证明：若\Term{functor}{函子} $h$ 与\Term{filtered-direct-limit}{滤过归纳极限}交换，则在由满足 $u\in h(X')$ 的对 $(X',u)$ 组成的\Term{filtered}{滤过}\Term{category}{范畴} $E'_{/h}$ 中，由满足 $X'\in\operatorname{Ob}E'_{\mathrm{PF}}$ 的对 $(X',u)$ 组成的子\Term{category}{范畴}\Term{cofinal}{余尾}。亦即对 $E'_{/h}$ 的任意 $(X',u)$，存在\Term{morphism}{态射} $X'_0\to X'$，$X'_0\in\operatorname{Ob}E'_{\mathrm{PF}}$，使 $u\in\operatorname{Im}(h(X'_0)\to h(X'))$。已知 $X'$ 是 $E'_{\mathrm{PF}}$ 的\Term{object}{对象} $X'_i$ 的\Term{filtered-direct-limit}{滤过归纳极限}（1.24.2 a)），所以 $h(X')$ 是 $h(X'_i)$ 的\Term{filtered-direct-limit}{滤过归纳极限}，故得结论。
\end{proof}

\SourcePage{171}
f) 只须再证明蕴涵
\[
\text{(ii bis)}\Rightarrow\text{(ii)}\qquad\text{及}\qquad
\text{(iii)}\Leftrightarrow\text{(iii bis)}.
\]
第一个蕴涵显然由 $E'$ 的每个\Term{object}{对象} $X'$ 都是 $E'_{\mathrm{PF}}$ 的\Term{object}{对象} $X'_i$ 的\Term{filtered-direct-limit}{滤过归纳极限}得到。对第二个蕴涵，只须注意\Term{functor}{函子}族
\[
Y''\longmapsto\operatorname{Hom}(X'',Y''),\qquad X''\in\operatorname{Ob}E''_{\mathrm{coh}},
\]
\Term{conservative}{保守}，并且由与\Term{filtered-direct-limit}{滤过归纳极限}交换的\Term{functor}{函子}组成（1.23(ii)）；所以\Term{functor}{函子} $f:E'\to E''$ 与\Term{filtered-direct-limit}{滤过归纳极限}交换，当且仅当复合\Term{functor}{函子}
\[
X'\longmapsto\operatorname{Hom}(X',f(X'))
=\operatorname{Hom}_{\operatorname{Pro}(E')}(g_0(X''),X')
\]
如此。由4.12.1，这\Term{equivalence}{等价}于 $g_0(X'')$ 属于 $\operatorname{Pro}(X'_{\mathrm{PF}})$ 的\Term{essential-image}{本质像}。证毕。

\begin{SGAresult}{注记}{4.13}
宜用\Term{sheaf}{层}的语言说明条件4.10(iv)的含义。考虑\Term{functor}{函子}的\Term{commutative-diagram}{交换图}
\begin{equation}
\begin{tikzcd}
\operatorname{Pro}(E'_{\mathrm{coh}})^{\circ}\arrow[r,hook,"\alpha"]\arrow[d,"\beta_0"']&
\operatorname{Pro}(E')^{\circ}_{\beta}\arrow[d,"\beta"]\\
\check E'_{\mathrm{coh}}&\check E'\arrow[l,"\gamma"]
\end{tikzcd}\tag{4.13.1}
\end{equation}
这里 $\vee$ 表示取值于 $(\mathrm{Ens})$ 的\Term{functor}{函子}\Term{category}{范畴}，第二个水平\Term{arrow}{箭头} $\gamma$ 是\Term{functor}{函子}的\Term{restriction}{限制}\Term{arrow}{箭头}，其余\Term{arrow}{箭头}均为熟知的\Term{fully-faithful}{全忠实}\Term{functor}{函子}。注意，包含\Term{arrow}{箭头} $\alpha$ 一般不与\Term{inverse-limit}{射影极限}交换，即使\Term{finite}{有限}\Term{inverse-limit}{射影极限}也不例外；所以一般而言，\Term{site}{景} $C$（如 $E''$）上的\Term{sheaf}{层} $F$ 给出该\Term{site}{景}上取值于 $\operatorname{Pro}(E'')^{\circ}$ 的\Term{presheaf}{预层} $\alpha\circ F$，而它不一定是\Term{sheaf}{层}。不过，若 $C$ 上取值于 $\operatorname{Pro}(E'_{\mathrm{coh}})^{\circ}$ 的\Term{presheaf}{预层} $F$，其相应的取值于 $\operatorname{Pro}(E')^{\circ}$ 的\Term{presheaf}{预层} $\alpha F$ 是\Term{sheaf}{层}，则 $F$ 本身是\Term{sheaf}{层}。这是由于\Term{sheaf}{层}条件是一种\Term{left-exact}{左正合}性性质（第II讲6.2 1)），而\Term{functor}{函子} $\beta,\beta_0,\gamma$ 与小\Term{inverse-limit}{射影极限}交换。

另一方面，由4.9.10中关于 $E''_{\mathrm{coh}}$ 上取值于任意\Term{category}{范畴} $P$ 的\Term{sheaf}{层}的刻画，对于每个满足 $\alpha^{\circ}:P^{\circ}\to Q^{\circ}$ 与\Term{finite}{有限}和及按\Term{equivalence-relation}{等价关系}取商交换的\Term{functor}{函子} $\alpha:P\to Q$，从 $\operatorname{Hom}(E''_{\mathrm{coh}}{}^{\circ},P)$ 到 $\operatorname{Hom}(E''_{\mathrm{coh}}{}^{\circ},Q)$ 的\Term{functor}{函子} $g\mapsto\alpha\circ g'$ 把\Term{sheaf}{层}变为\Term{sheaf}{层}。另一方面，已知在 $\operatorname{Hom}E'_{\mathrm{coh}}$ 及在 $E'$ 中，\Term{finite}{有限}和及按\Term{equivalence-relation}{等价关系}取商\Term{representable}{可表}，且包含\Term{functor}{函子} $E'_{\mathrm{coh}}\to E'$ 与之交换（1.15及1.17.1）。
\SourcePage{172}
所以，对于\Term{category}{范畴} $\operatorname{Pro}(E'_{\mathrm{coh}})$、$\operatorname{Pro}(E')$ 及它们之间的包含\Term{functor}{函子}也如此（第I讲8.9.5 b)及8.9.7）；该\Term{functor}{函子}正是 $\alpha^{\circ}$，其中 $\alpha$ 是(4.13.1)中的\Term{functor}{函子}。由此得到：若
\[
\varphi:E'_{\mathrm{coh}}{}^{\circ}\longrightarrow\operatorname{Pro}(E'_{\mathrm{coh}})^{\circ}\tag{4.13.2}
\]
是\Term{functor}{函子}，则 $\varphi$ 是 $E''_{\mathrm{coh}}$ 上取值于 $\operatorname{Pro}(E'_{\mathrm{coh}})^{\circ}$ 的\Term{sheaf}{层}，当且仅当 $\alpha\circ\varphi$ 是 $E''_{\mathrm{coh}}$ 上取值于 $\operatorname{Pro}(E')^{\circ}$ 的\Term{sheaf}{层}。所以，满足条件4.10(iv)的\Term{gluing-functor}{粘合函子} $f:E'\to E''$ 的\Term{category}{范畴}，\Term{canonical}{典范}\Term{equivalence}{等价}于 $E''_{\mathrm{coh}}$ 上取值于 $\operatorname{Pro}(E'_{\mathrm{coh}})^{\circ}$ 的\Term{sheaf}{层} $\varphi$ 的\Term{category}{范畴}。结合4.10得知，这样的\Term{sheaf}{层} $\varphi$ \Term{canonical}{典范}地定义一个\Term{gluing-functor}{粘合函子}（准确到唯一\Term{isomorphism}{同构}），而由后者得到的\Term{topos}{拓扑斯} $E$ \Term{coherent}{凝聚}；此外，若 $E'$ \Term{perfect}{完美}，由4.11，这样得到所有能由\Term{topos}{拓扑斯} $E',E''$ \Term{gluing}{粘合}得到的\Term{coherent}{凝聚}\Term{topos}{拓扑斯}。

另一方面，注意 $E'_{\mathrm{coh}}$ \Term{equivalence}{等价}于小\Term{category}{范畴}，并对\Term{finite}{有限}\Term{inverse-limit}{射影极限}\Term{stable}{稳定}，所以由(4.13.1)中的\Term{fully-faithful}{全忠实}\Term{functor}{函子} $\beta_0$，$\operatorname{Pro}(E'_{\mathrm{coh}})^{\circ}$ \Term{equivalence}{等价}于 $\check E'_{\mathrm{coh}}=\operatorname{Hom}(E'_{\mathrm{coh}},(\mathrm{Ens}))$ 中由\Term{left-exact}{左正合}\Term{functor}{函子}组成的\Term{full-subcategory}{满子范畴}（第I讲8.10.14）。因此，给定\Term{sheaf}{层}(4.13.2)，\Term{equivalence}{等价}于给定\Term{functor}{函子}
\[
F:E''_{\mathrm{coh}}{}^{\circ}\times E'_{\mathrm{coh}}\longrightarrow(\mathrm{Ens}),\tag{4.13.3}
\]
它对第一个变元是\Term{sheaf}{层}，对第二个变元\Term{left-exact}{左正合}；亦即给定\Term{left-exact}{左正合}\Term{functor}{函子}
\[
f_0:E'_{\mathrm{coh}}\longrightarrow E''.\tag{4.13.4}
\]
若 $f$ 是与 $\varphi$ 关联且满足条件4.10(iv)的\Term{gluing-functor}{粘合函子}，则立即验证，$f_0$ \Term{canonical}{典范}\Term{isomorphism}{同构}于\Term{functor}{函子} $f:E'\to E''$ 的\Term{restriction}{限制}。所以在所考虑的条件下，$f_0$ 确定 $f$，准确到\Term{canonical}{典范}\Term{isomorphism}{同构}。最后，概括4.11的内容，指出：
\end{SGAresult}

\begin{SGAresult}{推论}{4.14}
当 $E'$ 是\Term{perfect}{完美}\Term{topos}{拓扑斯}时，\Term{gluing}{粘合}后得到\Term{coherent}{凝聚}\Term{topos}{拓扑斯} $E$ 的\Term{gluing-functor}{粘合函子} $f:E'\to E''$ 的\Term{category}{范畴}，通过 $f\mapsto f_0=f_{/E'_{\mathrm{coh}}}$，\Term{equivalence}{等价}于\Term{left-exact}{左正合}\Term{functor}{函子} $f_0:E'_{\mathrm{coh}}\to E''$ 的\Term{category}{范畴}；亦即\Term{equivalence}{等价}于\Term{site}{景} $E''_{\mathrm{coh}}$（\Term{equivalence}{等价}地，\Term{topos}{拓扑斯} $E''$）上取值于\Term{category}{范畴} $\operatorname{Pro}(E'_{\mathrm{coh}})^{\circ}$ 的\Term{sheaf}{层} $\varphi$ 的\Term{category}{范畴}（见4.9.10）。
\end{SGAresult}
注意，第一个断言也可由4.10得到：只须注意 $E'\simeq\operatorname{Ind}(E'_{\mathrm{coh}})$（1.25），所以，与\Term{filtered-direct-limit}{滤过归纳极限}交换的\Term{functor}{函子} $f:E'\to P$ 的\Term{category}{范畴}，通过\Term{restriction}{限制}\Term{functor}{函子}\Term{equivalence}{等价}于任意\Term{functor}{函子} $f:E'_{\mathrm{coh}}\to P$ 的\Term{category}{范畴}；这里 $P$ 是小\Term{filtered-direct-limit}{滤过归纳极限}\Term{representable}{可表}的\Term{category}{范畴}。立即看出，在这些条件下，若 $P$ 中的\Term{finite}{有限}\Term{direct-limit}{归纳极限}\Term{representable}{可表}，则 $f$ \Term{left-exact}{左正合}，当且仅当 $f_0$ \Term{left-exact}{左正合}。

\SourcePage{173}
\begin{SGAresult}{注记}{4.15}
a) 4.10的证明表明：以一个定义\Term{conservative}{保守}\Term{fiber-functor}{纤维函子}族的子\Term{category}{范畴}代替\Term{category}{范畴} $\operatorname{Point}(E'')$，便得到与(ii ter)（相应地，(iii ter)）\Term{equivalence}{等价}的条件。

b) 取 $E''$ 为\Term{point-topos}{点拓扑斯}（第IV讲2.2），立即看出条件(ii ter)与(iii ter)\Term{equivalence}{等价}，只有当 $E'_{\mathrm{coh}}=E'_{\mathrm{PF}}$，即 $E'$ \Term{perfect}{完美}时才成立。因此，条件“$E$ \Term{coherent}{凝聚}”与条件“$f$ 与\Term{filtered-direct-limit}{滤过归纳极限}交换”\Term{equivalence}{等价}，只有当 $E'$ \Term{perfect}{完美}时才成立。反过来，更一般地取 $E''$ 为形如 $\widehat C$ 的\Term{coherent}{凝聚}\Term{topos}{拓扑斯}，则在这种情况下，对任意 $E'$，条件(i)都\Term{equivalence}{等价}于(iv)。事实上，撰稿人尚未能在一般情形下构造出一个 $E$ \Term{coherent}{凝聚}而不满足条件(iv)的例子。这真令人羞愧！
\end{SGAresult}

\SGAsection{5}{函子 $H^i(X,-)$ 与滤过归纳极限的交换性}
\begin{SGAresult}{定理}{5.1}
设 $E'$ 是\Term{algebraic}{代数}\Term{topos}{拓扑斯}（2.3），$E$ 是\Term{locally-coherent}{局部凝聚}\Term{topos}{拓扑斯}（2.3），$f:E'\to E$ 是\Term{coherent}{凝聚}\Term{topos-morphism}{拓扑斯态射}（3.1）。则对于每个整数 $q$，\Term{functor}{函子} $R^qf_*$（第V讲5.0）与\Term{abelian-sheaf}{阿贝尔层}的\Term{filtered-direct-limit}{滤过归纳极限}交换。
\end{SGAresult}

\begin{SGAresult}{推论}{5.2}
设 $E'$ 是\Term{coherent}{凝聚}\Term{topos}{拓扑斯}（2.3）。对每个整数 $q$，\Term{functor}{函子} $H^q(E',-)$ 与\Term{abelian-sheaf}{阿贝尔层}的\Term{filtered-direct-limit}{滤过归纳极限}交换。
\end{SGAresult}

\begin{SGAresult}{推论}{5.3}
设 $E$ 是\Term{topos}{拓扑斯}（相应地，\Term{algebraic}{代数}\Term{topos}{拓扑斯}）（2.3），$X$ 是 $E$ 的\Term{algebraic}{代数}\Term{coherent}{凝聚}\Term{object}{对象}（相应地，\Term{coherent}{凝聚}\Term{object}{对象}）（2.3）。对每个整数 $q$，\Term{functor}{函子} $H^q(X,-)$ 与\Term{abelian-sheaf}{阿贝尔层}的\Term{filtered-direct-limit}{滤过归纳极限}交换。
\end{SGAresult}

取 $E$ 为\Term{point-topos}{点拓扑斯}，$f:E'\to E$ 为唯一的\Term{morphism}{态射}（第IV讲2.2），则推论5.2由5.1得到。来证明5.2蕴涵5.3。设 $j_X:E_{/X}\to E$ 为\Term{localization}{局部化}\Term{morphism}{态射}。有关于\Term{abelian-sheaf}{阿贝尔层}\Term{functor}{函子}性的\Term{canonical}{典范}\Term{isomorphism}{同构}（第V讲2.2.1）
\[
H^q(X,F)\simeq H^q(E_{/X},j_X^*F).
\]
\Term{functor}{函子} $j_X^*$ 与\Term{direct-limit}{归纳极限}交换，且\Term{topos}{拓扑斯} $E_{/X}$ \Term{coherent}{凝聚}（2.4.6），所以得到5.3。来证明5.3蕴涵5.1。设 $E_{\mathrm{cohalg}}$ 为 $E$ 中由\Term{algebraic}{代数}\Term{coherent}{凝聚}\Term{object}{对象}组成的\Term{full-subcategory}{满子范畴}。它是生成\Term{category}{范畴}（2.4.5）。对每个\Term{abelian-sheaf}{阿贝尔层} $F$，$R^qf_*F$ 是\Term{presheaf}{预层} $X\mapsto H^q(f^*X,F)$（$X\in\operatorname{Ob}C$）的\Term{associated-sheaf}{关联层}（第V讲5.1）。由于 $f$ \Term{coherent}{凝聚}，对 $C$ 的每个\Term{object}{对象} $X$，$f^*X$ \Term{coherent}{凝聚}（3.1）。由5.3，\Term{presheaf}{预层} $X\mapsto H^q(f^*X,F)$ 与变元 $F$ 的\Term{filtered-direct-limit}{滤过归纳极限}交换；利用\Term{associated-sheaf}{关联层}\Term{functor}{函子}与\Term{direct-limit}{归纳极限}交换，得到5.1。（注意，我们只用了\Term{morphism}{态射} $f$ 的下述性质：存在 $E$ 的\Term{generating-family}{生成族} $S$，使对 $S$ 的每个 $X$，$f^*(X)$ 是\Term{algebraic}{代数}且\Term{coherent}{凝聚}的。）所以只须证明5.2。已知\Term{functor}{函子} $H^0(E',-)$ 与\Term{abelian-sheaf}{阿贝尔层}的\Term{filtered-direct-limit}{滤过归纳极限}交换（1.2.3）。
\SourcePage{174}
由[7]中的经典结果，为证5.2，只须证明：对于\Term{functor}{函子} $H^0(E',-)$ \Term{acyclic}{无上同调}的\Term{sheaf}{层}的\Term{filtered-direct-limit}{滤过归纳极限}，仍对 $H^0(E',-)$ \Term{acyclic}{无上同调}。这由下述引理应用于 $C=E'_{\mathrm{coh}}$ 得到。

\begin{SGAresult}{引理}{5.4}
设 $E$ 是\Term{locally-coherent}{局部凝聚}\Term{topos}{拓扑斯}（2.3），$C$ 是 $E$ 的\Term{full-subcategory}{满子范畴}，它生成、对\Term{fiber-product}{纤维积}\Term{stable}{稳定}，且其每个\Term{object}{对象}都\Term{quasi-compact}{拟紧}（1.1）。则 $C$-\Term{acyclic}{无上同调}\Term{sheaf}{层}（第V讲4.2）的\Term{filtered-direct-limit}{滤过归纳极限}仍是 $C$-\Term{acyclic}{无上同调}\Term{sheaf}{层}。
\end{SGAresult}
\begin{proof}
设 $(F_i)_{i\in I}$ 是 $C$-\Term{acyclic}{无上同调}\Term{sheaf}{层}的\Term{filtered}{滤过}\Term{inductive-system}{归纳系统}，以 $F$ 表示其\Term{direct-limit}{归纳极限}。$C$ 的每个\Term{object}{对象} $Y$ 都\Term{coherent}{凝聚}（2.1），从而
\[
F(Y)=\varinjlim_i F_i(Y)\qquad\text{（1.2.3）。}
\]
为证 $F$ 是 $C$-\Term{acyclic}{无上同调}的，只须证明对每个 $q>0$ 和每个 $Y\in\operatorname{Ob}C$，都有 $\check H^q(Y,F)=0$（第V讲4.3）。设 $X$ 是 $C$ 的\Term{object}{对象}，$\mathscr X=(X_\alpha\to X)_{\alpha\in A}$ 是由 $C$ 的\Term{object}{对象}组成的\Term{finite}{有限}\Term{covering-family}{覆盖族}。有
\[
H^q(\mathscr X,F)=H^q(C^{\bullet}(\mathscr X,F))\qquad\text{（第V讲2.4.3）。}
\]
由前述结果，并利用 $C^{\bullet}(\mathscr X,F)$ 只涉及群的\Term{finite-product}{有限积}（第V讲2.3.3），有
\[
C^{\bullet}(\mathscr X,F)=\varinjlim_i C^{\bullet}(\mathscr X,F_i).
\]
所以对于 $q>0$，
\[
\begin{aligned}
H^q(\mathscr X,F)
&=H^q(C^{\bullet}(\mathscr X,F))
=H^q\bigl(\varinjlim_i C^{\bullet}(\mathscr X,F_i)\bigr)\\
&=\varinjlim_i C^{\bullet}(\mathscr X,F_i)=0
\end{aligned}
\]
（第V讲4.3）。因此 $\check H^q(X,F)=\varinjlim_{\mathscr X}H^q(\mathscr X,F)=0$。
\end{proof}

\begin{SGAresult}{推论}{5.5}
设 $E$ 是\Term{coherent}{凝聚}\Term{topos}{拓扑斯}（2.3），$(X_i)_{i\in I}$ 是其\Term{final-object}{终对象}的\Term{coherent}{凝聚}\Term{subobject}{子对象}所成的递减\Term{filtered}{滤过}族。以 $\Phi$ 表示 $E$ 的那些包含在某个 $X_i$ 的补\Term{closed-object}{闭对象}中的\Term{closed-object}{闭对象}所成的族（第IV讲9）。则 $\Phi$ 是 $E$ 的\Term{support-family}{支撑族}（第V讲6.12），并且对于每个整数 $q$，\Term{functor}{函子} $H^q_\Phi(E,-)$ 及 $\mathscr H^q_\Phi$（第V讲6.13）与\Term{filtered-direct-limit}{滤过归纳极限}交换。
\end{SGAresult}
\begin{proof}
设 $Z_i$ 是 $X_i$ 的补\Term{closed-object}{闭对象}。对每个\Term{abelian-sheaf}{阿贝尔层} $F$，有\Term{exact-sequence}{正合列}（第V讲6.5）：
\[
0\longrightarrow H^0_{Z_i}(E,F)\longrightarrow H^0(E,F)
\longrightarrow H^0(X_i,F)\longrightarrow H^1_{Z_i}(E,F)\longrightarrow\cdots.
\]
\Term{functor}{函子} $H^q(E,F)$ 与 $H^q(X_i,F)$ 与变元 $F$ 的\Term{direct-limit}{归纳极限}交换（5.2、5.3）。所以 $H^q_{Z_i}(E,F)$ 与变元 $F$ 的\Term{filtered-direct-limit}{滤过归纳极限}交换；对 $Z_i$ 取\Term{direct-limit}{归纳极限}，便得到\Term{functor}{函子} $H^q_\Phi(E,F)$ 的相应性质（第V讲6.13）。关于\Term{functor}{函子} $\mathscr H^q_\Phi$ 的断言，则由\Term{localization}{局部化}并取\Term{associated-sheaf}{关联层}得到。
\end{proof}

\begin{SGAresult}{定义}{5.6}
设 $E$ 是以 $A$ 为环的\Term{ringed-topos}{环化拓扑斯}，$q$ 是整数。若存在\Term{covering}{覆盖} $E$ 的\Term{final-object}{终对象}的\Term{object}{对象}族 $X_i$，使对每个 $i$，在 $E_{/X_i}$ 上有模的\Term{exact-sequence}{正合列}
\[
L_q\longrightarrow L_{q-1}\longrightarrow\cdots L_1\longrightarrow L_0\longrightarrow G_{/X_i}\longrightarrow0,\tag{5.5.1}
\]
其中对于每个 $p$，$L_p=A_{/X_i}^{n_p}$，$n_p$ 为整数，则称 $A$-模 $G$ 为 $q$-\Term{finite-presentation}{有限表现}的。
\end{SGAresult}

\begin{SGAresult}{命题}{5.7}
设 $G$ 是 $q$-\Term{finite-presentation}{有限表现}的\Term{sheaf}{层}。对每个 $i\leq q-1$，\Term{functor}{函子} $\mathscr{E}\!xt_A^i(G,-)$ 与 $A$-模的\Term{filtered-direct-limit}{滤过归纳极限}交换。
\end{SGAresult}
\SourcePage{175}
\begin{proof}
此问题在 $E$ 上是\Term{local-scheme}{局部}的（第V讲6.1）。所以\Term{localization}{局部化}到 $E_{/X_i}$ 后，可以假定有形如(5.5.1)的\Term{exact-sequence}{正合列}。置 $L_{\bullet}=L_q\to\cdots L_1\to L_0$。有
\[
\mathscr{E}\!xt_A^i(G,F):\mathscr H^i(\mathscr{H}\!om_A(L_{\bullet},F)),
\]
而\Term{functor}{函子} $F\mapsto\mathscr{H}\!om_A(L_{\bullet},F)$ 与\Term{direct-limit}{归纳极限}交换，故得命题。
\end{proof}

\begin{SGAresult}{推论}{5.8}
设 $(E,A)$ 是\Term{ringed-topos}{环化拓扑斯}，$G$ 是 $q$-\Term{finite-presentation}{有限表现}的 $A$-模\Term{sheaf}{层}，$X$ 是 $E$ 的\Term{algebraic}{代数}\Term{coherent}{凝聚}\Term{object}{对象}。对满足 $i(i-1)/2\leq q-1$ 的 $i$，\Term{functor}{函子} $F\mapsto\operatorname{Ext}_A^i(X;G,F)$ 与变元 $F$ 的\Term{filtered-direct-limit}{滤过归纳极限}交换。
\end{SGAresult}
由5.7、5.3及\Term{spectral-sequence}{谱序列}（第V讲6.1.3）得到。

\textbf{5.9.} 设 $(E,A)$ 是\Term{ringed-topos}{环化拓扑斯}，$\Phi$ 是 $E$ 的\Term{support-family}{支撑族}（第V讲6.12），$G$ 是 $r$-\Term{finite-presentation}{有限表现}的 $A$-模（5.6），$X$ 是 $E$ 的\Term{algebraic}{代数}\Term{coherent}{凝聚}\Term{object}{对象}（2.3）。分别在第V讲6.9.1和第V讲6.9.2中的最后一个\Term{spectral-sequence}{谱序列}中，对 $\Phi$ 中的\Term{closed-object}{闭对象}取\Term{direct-limit}{归纳极限}，得到两个\Term{spectral-sequence}{谱序列}（第V讲6.13）：
\[
\begin{cases}
{}'E_2^{pq}=\varinjlim_{Z\in\Phi}\operatorname{Ext}_A^p(X;G,\mathscr H_Z^qF)
\Rightarrow\operatorname{Ext}_{A,\Phi}^{p+q}(X;G,F),\\[3pt]
{}''E_2^{pq}=\varinjlim_{Z\in\Phi}\mathscr{E}\!xt_A^p(G,\mathscr H_Z^qF)
\Rightarrow\mathscr{E}\!xt_{A,\Phi}^{p+q}(G,F).
\end{cases}\tag{5.9.1}
\]
由5.7及5.8，有\Term{canonical}{典范}\Term{isomorphism}{同构}
\[
\begin{cases}
{}'E_2^{pq}\simeq\operatorname{Ext}_A^p(X;G,\mathscr H_\Phi^qF),&\dfrac{p(p-1)}2\leq r-1,\\[5pt]
{}''E_2^{pq}\simeq\mathscr{E}\!xt_A^p(G,\mathscr H_\Phi^qF),&p\leq r-1.
\end{cases}\tag{5.9.2}
\]

\begin{SGAresult}{推论}{5.10}
采用5.5的假设和记号。设 $A$ 是 $E$ 的环，$G$ 是 $q$-\Term{finite-presentation}{有限表现}的 $A$-模。对每个满足
\[
\frac{i(i-1)}2\leq r-1
\]
的整数 $i$，\Term{functor}{函子} $F\mapsto\operatorname{Ext}_{A,\Phi}^i(E;G,F)$ 及 $F\mapsto\mathscr{E}\!xt_{A,\Phi}^i(G,F)$（第V讲6.13）与\Term{filtered-direct-limit}{滤过归纳极限}交换。
\end{SGAresult}
分别由5.5、5.7及第V讲6.9.1和第V讲6.9.2中的第一个\Term{spectral-sequence}{谱序列}得到。

\SGAsection{6}{纤维范畴的归纳极限与射影极限}
\textbf{6.0.} 假定读者熟悉\Term{fibered-category}{纤维范畴}理论[5]。本节主要是为了固定关于\Term{fibered-category}{纤维范畴}的术语和记号。除非另有说明，本节所考虑的所有\Term{category}{范畴}都属于一个固定的\Term{universe}{宇宙} $\mathcal U$。

设 $\mathscr F,\mathscr G,\mathscr E$ 是三个\Term{category}{范畴}，$\pi:\mathscr F\to\mathscr E$、$\pi':\mathscr G\to\mathscr E$ 是两个\Term{functor}{函子}。以 $\mathscr{H}\!om_{\mathscr E}(\mathscr F,\mathscr G)$ 表示如下\Term{category}{范畴}：其\Term{object}{对象}是满足图
\SourcePage{176}
\[
\begin{tikzcd}
\mathscr F\arrow[rr,"u"]\arrow[dr,"\pi"']&&\mathscr G\arrow[dl,"\pi'"]\\
&\mathscr E&
\end{tikzcd}
\]
交换的\Term{functor}{函子} $u:\mathscr F\to\mathscr G$；其\Term{morphism}{态射}是\Term{functor}{函子}的 $\mathscr E$-\Term{morphism}{态射}，即由 $\pi'$ 变成恒等\Term{morphism}{态射}的\Term{functor}{函子}\Term{morphism}{态射}。

设 $\xi$ 是底\Term{category}{范畴} $\mathscr E$ 的\Term{object}{对象}。满足 $\pi(X)=\xi$ 的 $\mathscr F$ 的\Term{object}{对象} $X$ 称为 $\xi$ 上方的 $\mathscr F$ 的\Term{object}{对象}。设 $f:\xi\to\eta$ 是 $\mathscr E$ 的\Term{morphism}{态射}。满足 $\pi(m)=f$ 的 $\mathscr F$ 的\Term{morphism}{态射} $m$ 称为 $f$ 上方的 $\mathscr F$ 的\Term{morphism}{态射}。设 $f:\xi\to\eta$ 是 $\mathscr E$ 的\Term{morphism}{态射}，$X$ 是 $\xi$ 上方的 $\mathscr F$ 的\Term{object}{对象}，$Y$ 是 $\eta$ 上方的 $\mathscr F$ 的\Term{object}{对象}；以 $\operatorname{Hom}_f(X,Y)$ 表示 $f$ 上方的从 $X$ 到 $Y$ 的\Term{morphism}{态射}集。

\begin{SGAresult}{定义}{6.1}
1) 若对于 $\pi(m)$ 上方的每个\Term{morphism}{态射} $p:Z\to Y$，存在唯一的 $\pi(X)$ 的恒等\Term{morphism}{态射}上方的\Term{morphism}{态射} $q:Z\to X$，使 $mq=p$，则称 $\mathscr F$ 的\Term{morphism}{态射} $m:X\to Y$ 为\Term{cartesian-morphism}{笛卡尔态射}（morphisme cartésien）。

2) 若对于 $\mathscr E$ 的每个\Term{morphism}{态射} $f:\xi\to\eta$，$\eta$ 上方的 $\mathscr F$ 的每个\Term{object}{对象}都是某个 $f$ 上方的\Term{cartesian-morphism}{笛卡尔态射}的目标，则称\Term{category}{范畴} $\mathscr F$ 由 $\pi$ 在 $\mathscr E$ 上\Term{prefibered}{预纤维化}（préfibrée）。

3) 若\Term{category}{范畴} $\mathscr F$ \Term{prefibered}{预纤维化}，且 $\mathscr F$ 中两个可复合\Term{cartesian-morphism}{笛卡尔态射}的复合仍是\Term{cartesian-morphism}{笛卡尔态射}，则称 $\mathscr F$ 由 $\pi$ 在 $\mathscr E$ 上\Term{fibered}{纤维化}。
\end{SGAresult}

\textbf{6.1.1.} 设 $\xi$ 是 $\mathscr E$ 的\Term{object}{对象}。$\mathscr F$ 在 $\xi$ 处的\Term{fibered-category}{纤维范畴}，记作 $\mathscr F_\xi$，是如下 $\mathscr F$ 的子\Term{category}{范畴}：其\Term{object}{对象}是 $\xi$ 上方的 $\mathscr F$ 的\Term{object}{对象}，\Term{morphism}{态射}是 $\operatorname{id}_\xi$ 上方的 $\mathscr F$ 的\Term{morphism}{态射}。设 $X,Y$ 是 $\mathscr F_\xi$ 的两个\Term{object}{对象}，以 $\operatorname{Hom}_\xi(X,Y)$ 表示 $\mathscr F_\xi$ 中从 $X$ 到 $Y$ 的\Term{morphism}{态射}集。

\textbf{6.1.2.} 设 $f:\xi\to\eta$ 是 $\mathscr E$ 的\Term{morphism}{态射}。$\eta$ 上方的\Term{object}{对象} $Y$ 是某个 $f$ 上方的\Term{cartesian-morphism}{笛卡尔态射}的目标，当且仅当定义在\Term{fiber}{纤维} $\mathscr F_\xi$ 上的\Term{set}{集合}值\Term{functor}{函子}
\[
Z\longmapsto\operatorname{Hom}_f(Z,Y),\qquad Z\in\operatorname{Ob}(\mathscr F_\xi),
\]
\Term{representable}{可表}。当\Term{category}{范畴} $\mathscr F$ \Term{prefibered}{预纤维化}时，取值于 $\mathscr F_\xi$ 上的\Term{set}{集合}\Term{presheaf}{预层}\Term{category}{范畴}的\Term{functor}{函子}
\[
Y\longmapsto\operatorname{Hom}_f(-,Y),\qquad Y\in\operatorname{Ob}(\mathscr F_\eta),
\]
实际上取值于 $\mathscr F_\xi$ 上的\Term{representable}{可表}\Term{functor}{函子}\Term{category}{范畴}，所以它确定一个\Term{functor}{函子} $f^*:\mathscr F_\eta\to\mathscr F_\xi$，准确到唯一\Term{isomorphism}{同构}；此\Term{functor}{函子}称为关于 $f$ 的\Term{inverse-image}{逆像}\Term{functor}{函子}。
\SourcePage{177}
假定\Term{category}{范畴} $\mathscr F$ \Term{prefibered}{预纤维化}，并对 $\mathscr E$ 的每个\Term{morphism}{态射} $f$ 选定\Term{base-change}{基变换}\Term{functor}{函子} $f^*$。则\Term{category}{范畴} $\mathscr F$ 在 $\mathscr E$ 上\Term{fibered}{纤维化}，当且仅当对 $\mathscr E$ 的每对可复合\Term{morphism}{态射} $f,g$，复合\Term{functor}{函子} $f^*g^*$ 都是\Term{inverse-image}{逆像}\Term{functor}{函子}。此时以
\[
C_{f,g}:f^*g^*\longrightarrow(gf)^*\tag{6.1.2.1}
\]
表示\Term{canonical}{典范}\Term{isomorphism}{同构}。由 $\mathscr E$ 中\Term{morphism}{态射}复合的结合律，\Term{isomorphism}{同构} $C_{f,g}$ 满足\Term{cocycle-condition}{余圈条件}：
\[
C_{f,hg}:(f^*\circ C_{g,h})=C_{gh,f}(C_{f,g}\circ h^*).\tag{6.1.2.2}
\]

\textbf{6.1.3.} 反之，若对 $\mathscr E$ 的每个\Term{object}{对象} $\xi$ 给定\Term{category}{范畴} $\mathscr F_\xi$，对每个\Term{morphism}{态射} $f:\xi\to\eta$ 给定\Term{functor}{函子} $f^*:\mathscr F_\eta\to\mathscr F_\xi$，并对 $\mathscr E$ 的每对\Term{morphism}{态射} $(f,g)$ 给定\Term{functor}{函子}\Term{isomorphism}{同构} $C_{f,g}:f^*g^*\to(fg)^*$，使各 $C_{f,g}$ 满足条件(6.1.2.2)，则可以以本质唯一的方式构造 $\mathscr E$ 上的\Term{fibered-category}{纤维范畴} $\mathscr F$，其\Term{fiber}{纤维}“就是”各\Term{category}{范畴} $\mathscr F_\xi$，其\Term{base-change}{基变换}\Term{functor}{函子}可以选为“等于”预先给定的\Term{functor}{函子}[5]。

\textbf{6.1.4.} 设 $\mathscr F,\mathscr G$ 是 $\mathscr E$ 上的两个\Term{category}{范畴}。以
\[
\mathscr{H}\!om_{\mathrm{cart}/\mathscr E}(\mathscr F,\mathscr G)
\]
表示 $\mathscr{H}\!om_{\mathscr E}(\mathscr F,\mathscr G)$ 中由把 $\mathscr F$ 的\Term{cartesian-morphism}{笛卡尔态射}变成 $\mathscr G$ 的\Term{cartesian-morphism}{笛卡尔态射}的\Term{functor}{函子}所确定的\Term{full-subcategory}{满子范畴}。$\mathscr{H}\!om_{\mathrm{cart}/\mathscr E}(\mathscr F,\mathscr G)$ 的\Term{object}{对象}称为\Term{cartesian-functor}{笛卡尔函子}。

设 $C,C'$ 是两个\Term{category}{范畴}，$S$ 是\Term{category}{范畴} $C$ 的\Term{morphism}{态射}集。以 $\operatorname{Hom}_{S^{-1}}(C,C')$ 表示把 $S$ 的\Term{morphism}{态射}变成\Term{isomorphism}{同构}的从 $C$ 到 $C'$ 的\Term{functor}{函子}集。

以 $(\mathrm{Cat})$ 表示如下\Term{category}{范畴}：其\Term{object}{对象}是属于该\Term{universe}{宇宙}的\Term{category}{范畴}，\Term{morphism}{态射}是这些\Term{category}{范畴}之间的\Term{functor}{函子}（\Term{category}{范畴} $(\mathrm{Cat})$ 不属于该\Term{universe}{宇宙}）。

设 $\mathscr F$ 是 $\mathscr E$ 上的\Term{category}{范畴}，$S$ 是 $\mathscr F$ 的\Term{cartesian-morphism}{笛卡尔态射}集。\Term{category}{范畴} $\mathscr F$ 在 $(\mathrm{Cat})$ 上定义两个取值于属于该\Term{universe}{宇宙}的\Term{set}{集合}所成\Term{category}{范畴}的\Term{functor}{函子}：
\[
C\longmapsto\operatorname{Hom}_{S^{-1}}(\mathscr F,C),\qquad C\in\operatorname{Ob}(\mathrm{Cat}),
\]
\[
\begin{aligned}
C\longmapsto\operatorname{Hom}_{\mathrm{Cart}/\mathscr E}(\mathscr F,C\times\mathscr E)
=\{\text{从 $\mathscr F$ 到 $C\times\mathscr E$ 的笛卡尔函子}\}.
\end{aligned}
\]
（通过第二\Term{projection}{投影}\Term{functor}{函子}，把\Term{category}{范畴} $C\times\mathscr E$ 看作 $\mathscr E$ 上的\Term{category}{范畴}。）

\SourcePage{178}
\begin{SGAresult}{命题}{6.2}
\Term{functor}{函子} $(\mathrm{Cat})\to(\mathrm{Ens})$
\[
C\longmapsto\operatorname{Hom}_{S^{-1}}(\mathscr F,C),\qquad
C\longmapsto\operatorname{Hom}_{\mathrm{Cart}/\mathscr E}(\mathscr F,C\times\mathscr E)
\]
\Term{canonical}{典范}\Term{isomorphism}{同构}，并且\Term{representable}{可表}。
\end{SGAresult}
\begin{proof}
为证第一个断言，只须注意 $C\times\mathscr E$ 的\Term{cartesian-morphism}{笛卡尔态射}是形如 $m\times f$、其中 $m$ 为 $C$ 的\Term{isomorphism}{同构}的\Term{morphism}{态射}。为证第二个断言，只须证明\Term{functor}{函子} $\operatorname{Hom}_{S^{-1}}(\mathscr F,-)$ \Term{representable}{可表}。这由[1]得到。这里只指出证明的思路。形式地将 $S$ 的\Term{morphism}{态射}的逆添加到 $\mathscr F$ 的\Term{morphism}{态射}中。考虑由 $\mathscr F$ 的\Term{object}{对象}、$\mathscr F$ 的\Term{morphism}{态射}及 $S$ 的\Term{morphism}{态射}的形式逆生成的\Term{free-category}{自由范畴}（\Term{path-category}{道路范畴}）。按 $\mathscr F$ 中\Term{morphism}{态射}之间的关系以及下列形式的关系取商：
\[
s^{-1}s=\operatorname{id},\qquad ss^{-1}=\operatorname{id},\qquad s\in S.
\]
将所得\Term{category}{范畴}记作 $\mathscr F(S^{-1})$；它连同\Term{canonical}{典范}\Term{functor}{函子} $Q=\mathscr F\to\mathscr F(S^{-1})$ 表示\Term{functor}{函子} $\operatorname{Hom}_{S^{-1}}(\mathscr F,-)$。
\end{proof}

\begin{SGAresult}{定义}{6.3}
当 $\mathscr F$ 在 $\mathscr E$ 上\Term{fibered}{纤维化}时，将\Term{functor}{函子} $\operatorname{Hom}_{S^{-1}}(\mathscr F,-)$ 记作 $\varinjlim_{\mathscr E^{\circ}}\mathscr F$，称为 $\mathscr F$ 在 $\mathscr E^{\circ}$ 上的\Term{direct-limit}{归纳极限}\Term{functor}{函子}。表示它的\Term{category}{范畴}仍记作 $\varinjlim_{\mathscr E^{\circ}}\mathscr F$，称为 $\mathscr F$ 在 $\mathscr E^{\circ}$ 上的\Term{direct-limit}{归纳极限}\Term{category}{范畴}。
\end{SGAresult}

\textbf{6.4.0.} 假定 $\mathscr F$ 在 $\mathscr E$ 上\Term{fibered}{纤维化}；对 $\mathscr E$ 的每个\Term{morphism}{态射} $f$，选定 $f$ 上方的\Term{base-change}{基变换}\Term{functor}{函子}，并设
\[
Q:\mathscr F\longrightarrow\varinjlim_{\mathscr E^{\circ}}\mathscr F
\]
是\Term{canonical}{典范}\Term{functor}{函子}。$\mathscr F$ 的\Term{fiber}{纤维}到 $\mathscr F$ 的包含\Term{functor}{函子}与\Term{functor}{函子} $Q$ 复合，对 $\mathscr E$ 的每个\Term{object}{对象} $\xi$ 给出\Term{functor}{函子}
\[
u_\xi:\mathscr F_\xi\longrightarrow\varinjlim_{\mathscr E^{\circ}}\mathscr F,
\]
\SourcePage{179}
并对 $\mathscr E$ 的每个\Term{morphism}{态射} $f:\xi\to\eta$ 给出准确到\Term{canonical}{典范}\Term{isomorphism}{同构}交换的图：
\[
\begin{tikzcd}
\mathscr F_\eta\arrow[d,"f^*"']\arrow[dr,"u_\eta"]&\\
\mathscr F_\xi\arrow[r,"u_\xi"']&\varinjlim_{\mathscr E^{\circ}}\mathscr F
\end{tikzcd}
\]
于是，\Term{category}{范畴} $\varinjlim_{\mathscr E^{\circ}}\mathscr F$ 表现为如下\Term{pseudofunctor}{伪函子} $\mathscr E^{\circ}\to\mathrm{Cat}$ 的\Term{pseudofunctor}{伪函子}意义下的\Term{direct-limit}{归纳极限}[5]：它给每个 $\xi\in\operatorname{Ob}\mathscr E$ 关联 $\mathscr F_\xi$，给每个 $f:\xi\to\eta$ 关联\Term{functor}{函子} $f^*:\mathscr F_\eta\to\mathscr F_\xi$。不过注意，即使 $\xi\mapsto\mathscr F_\xi$ 是真正的\Term{functor}{函子}，\Term{category}{范畴} $\varinjlim_{\mathscr E^{\circ}}\mathscr F$ 一般也不是\Term{functor}{函子} $\xi\mapsto\mathscr F_\xi$ 在第I讲2的意义下的\Term{direct-limit}{归纳极限}（见6.8）。

\begin{SGAresult}{命题}{6.4}
设 $\mathscr F$ 是 $\mathscr E$ 上的\Term{fibered-category}{纤维范畴}。假定\Term{category}{范畴} $\mathscr E$ 有下列性质：

L1) 每个图
\[
\begin{tikzcd}[column sep=large]
\bullet\arrow[dr]&\\
&\bullet\\
\bullet\arrow[ur]&
\end{tikzcd}
\]
都可\Term{embedding}{嵌入}\Term{commutative-diagram}{交换图}
\[
\begin{tikzcd}
&\bullet\arrow[dr]&\\
\bullet\arrow[ur]\arrow[dr]&&\bullet\\
&\bullet\arrow[ur]&
\end{tikzcd}
\]
中。

L2) 对每对\Term{morphism}{态射} $u,v:\bullet\rightrightarrows\bullet$，若存在\Term{morphism}{态射} $t$ 满足关系 $tu=tv$，则存在\Term{morphism}{态射} $w$ 满足关系 $uw=vw$。

（注意，若\Term{category}{范畴} $\mathscr E^{\circ}$ \Term{pseudo-filtered}{伪滤过}（第I讲2.7），则它有性质L1)、L2)。）于是 $\mathscr F$ 的\Term{cartesian-morphism}{笛卡尔态射}集 $S$ 有下列性质：

Fr1) \Term{set}{集合} $S$ 对复合\Term{stable}{稳定}，且\Term{isomorphism}{同构}属于 $S$。
\SourcePage{180}
Fr2) 每个图
\[
\begin{tikzcd}[column sep=large]
\bullet\arrow[dr]&\\
&\bullet\\
\bullet\arrow[ur,"s"]&
\end{tikzcd}
\]
其中 $s$ 属于 $S$，都可补成\Term{commutative-diagram}{交换图}
\[
\begin{tikzcd}
&\bullet\arrow[dr]&\\
\bullet\arrow[ur,"t"]\arrow[dr]&&\bullet\\
&\bullet\arrow[ur,"s"']&
\end{tikzcd}
\]
其中 $t$ 属于 $S$。

Fr3) 对每对\Term{morphism}{态射} $u,v:\bullet\rightrightarrows\bullet$，若存在\Term{morphism}{态射} $s\in S$ 满足关系 $su=sv$，则存在\Term{morphism}{态射} $t\in S$ 使 $ut=vt$。
\end{SGAresult}
\begin{proof}
留给读者作为习题。
\end{proof}

\begin{SGAresult}{命题}{6.5}
设 $\mathscr F$ 是\Term{category}{范畴}，$S$ 是 $\mathscr F$ 的一个满足性质Fr1)、Fr2)、Fr3)（6.4）的\Term{morphism}{态射}集。对 $\mathscr F$ 的每个\Term{object}{对象} $X$，以 $S(X)$ 表示 $S$ 中以 $X$ 为目标的\Term{morphism}{态射}的\Term{category}{范畴}。\Term{category}{范畴} $S(X)$ \Term{cofiltered}{余滤过}（第I讲2.7）。\Term{category}{范畴} $\mathscr F(S^{-1})$（它表示\Term{functor}{函子} $\operatorname{Hom}_{S^{-1}}(\mathscr F,-)$）可如下描述：

a) $\mathscr F(S^{-1})$ 的\Term{object}{对象}就是 $\mathscr F$ 的\Term{object}{对象}。

b) 设 $X,Y$ 是 $\mathscr F(S^{-1})$ 的两个\Term{object}{对象}，则
\[
\operatorname{Hom}_{\mathscr F(S^{-1})}(X,Y)
=\varinjlim_{S(X)}\operatorname{Hom}_{\mathscr F}(-,Y).
\]

c) 设 $f:X\to Y$、$g:Y\to Z$ 是 $\mathscr F(S^{-1})$ 的两个\Term{morphism}{态射}。设
\[
\begin{tikzcd}
\bullet\arrow[d,"t"']\arrow[dr,"f'"]&\\
X&Y
\end{tikzcd}\qquad t\in S
\]
（相应地，
\[
\begin{tikzcd}
\bullet\arrow[d,"t'"']\arrow[dr,"g'"]&\\
Y&Z
\end{tikzcd}\qquad t\in S
\]
）为 $\mathscr F$ 中其像为 $f$（相应地，$g$）的图。
\SourcePage{181}
设
\[
\begin{tikzcd}
\bullet\arrow[d,"t''"']\arrow[r,"f''"]&\bullet\arrow[d,"t'"]\\
\bullet\arrow[r,"f'"']&Y
\end{tikzcd}\qquad t''\in S
\]
为\Term{commutative-diagram}{交换图}，其存在性由Fr2)保证。图
\[
\begin{tikzcd}
\bullet\arrow[d,"tt''"']\arrow[dr,"g'f''"]&\\
X&Z
\end{tikzcd}
\]
在 $\operatorname{Hom}_{\mathscr F(S^{-1})}(X,Z)$ 中的像就是复合\Term{morphism}{态射} $gf$。

设
\[
Q:\mathscr F\longrightarrow\mathscr F(S^{-1})
\]
为\Term{canonical}{典范}\Term{functor}{函子}。

d) \Term{functor}{函子} $Q^*:\widehat{\mathscr F(S^{-1})}\to\widehat{\mathscr F}$（$F\mapsto F\circ Q$，第I讲5.0）\Term{fully-faithful}{全忠实}，且在\Term{object}{对象}上\Term{injectivity}{单射}。

e) \Term{functor}{函子} $Q_!:\widehat{\mathscr F}\to\widehat{\mathscr F(S^{-1})}$（\Term{functor}{函子} $Q^*$ 的\Term{left-adjoint}{左伴随}，第I讲5.1）\Term{left-exact}{左正合}。（因此，当 $\mathscr F$ 中的\Term{finite}{有限}\Term{inverse-limit}{射影极限}\Term{representable}{可表}时，\Term{functor}{函子} $Q$ 也\Term{left-exact}{左正合}。）
\end{SGAresult}
\begin{proof}
证明参见[1]。
\end{proof}

\begin{SGAresult}{习题}{6.6}
设 $\mathscr F$ 是\Term{category}{范畴}，$S$ 是 $\mathscr F$ 的一个满足性质Fr1)、Fr2)、Fr3)的\Term{morphism}{态射}集。对 $\mathscr F$ 的每个\Term{object}{对象} $X$，以 $S(X)$ 表示以 $X$ 为目标的\Term{morphism}{态射}\Term{category}{范畴}。

a) \Term{category}{范畴} $S(X)$ \Term{cofiltered}{余滤过}（第I讲2）。

b) 对 $\mathscr F$ 的每个\Term{object}{对象} $X$，以 $J_S(X)$ 表示 $X$ 的那些包含由 $S$ 的一个\Term{morphism}{态射}生成的\Term{sieve}{筛}的\Term{sieve}{筛}所成的\Term{set}{集合}。各 $J_S(X)$ 在 $\mathscr F$ 上定义一个\Term{topology}{拓扑} $T$。

c) 关于\Term{topology}{拓扑} $T$，$S$ 的\Term{morphism}{态射}都是\Term{bicovering}{双覆盖}的（第II讲5.2）。

d) 对 $\mathscr F$ 的每个\Term{object}{对象} $X$，以 $\mathscr J^2(X)$ 表示以 $X$ 为目标的\Term{bicovering}{双覆盖}\Term{morphism}{态射}\Term{category}{范畴}。\Term{category}{范畴} $S(X)$ 在 $\mathscr J^2(X)$ 中\Term{cofinal}{余尾}（第I讲8）。

e) 设 $G$ 是 $\mathscr F$ 上的\Term{set}{集合}\Term{presheaf}{预层}。考虑关于\Term{topology}{拓扑} $T$ 的关联\Term{functor}{函子}。对 $\mathscr F$ 的每个\Term{object}{对象} $X$，有\Term{canonical}{典范}\Term{isomorphism}{同构}
\[
G(X)\xrightarrow{\sim}\lim_{S(X)}G(-).
\]
\SourcePage{182}
一个\Term{presheaf}{预层}是\Term{sheaf}{层}，当且仅当它把 $S$ 的每个\Term{morphism}{态射}变成\Term{isomorphism}{同构}。

f) 选择\Term{associated-sheaf}{关联层}\Term{functor}{函子}，使其在 $\mathscr F$ 上的\Term{restriction}{限制}在\Term{object}{对象}上\Term{injectivity}{单射}。以 $\mathscr F_{\to}$ 表示 $\mathscr F$ 上的\Term{set}{集合}\Term{sheaf}{层}\Term{category}{范畴}中由\Term{representable}{可表}\Term{presheaf}{预层}的\Term{associated-sheaf}{关联层}所确定的\Term{full-subcategory}{满子范畴}。\Term{category}{范畴} $\mathscr F_{\to}$ \Term{isomorphism}{同构}于6.5中描述的\Term{category}{范畴} $\mathscr F(S^{-1})$。$\mathscr F_{\to}$ 的\Term{object}{对象}构成 $\mathscr F$ 上的\Term{set}{集合}\Term{sheaf}{层}\Term{topos}{拓扑斯} $\mathscr F^{\sim}$ 的\Term{family-of-generators}{生成元族}。由 $\mathscr F^{\sim}$ 的\Term{canonical-topology}{典范拓扑}在 $\mathscr F_{\to}$ 上诱导的\Term{topology}{拓扑}是\Term{coarse-topology}{粗拓扑}（第II讲1.1.4）。

g) 设 $Q:\mathscr F\to\mathscr F_{\to}$ 为\Term{canonical}{典范}\Term{functor}{函子}。\Term{functor}{函子} $Q$ 是从\Term{site}{景} $\mathscr F_{\to}$（\Term{coarse-topology}{粗拓扑}）到\Term{site}{景} $\mathscr F$（b)中的\Term{topology}{拓扑} $T$）的\Term{site-morphism}{景态射}（第IV讲5.9），并诱导相应\Term{topos}{拓扑斯}之间的\Term{isomorphism}{同构}。

h) 设 $u:\mathscr F\to C$ 为\Term{functor}{函子}。\Term{functor}{函子} $u$ \Term{continuous}{连续}，当且仅当它把 $S$ 的\Term{morphism}{态射}变成\Term{isomorphism}{同构}（第III讲1.1）。

i) 设 $u:\mathscr F\to C$ 是把 $S$ 的\Term{morphism}{态射}变成\Term{isomorphism}{同构}的\Term{functor}{函子}。\Term{functor}{函子} $u$ 唯一地\Term{resolution}{分解}为
\[
\begin{tikzcd}
\mathscr F\arrow[rr,"u"]\arrow[dr,"Q"']&&C\\
&\mathscr F_{\to}\arrow[ur,"v"']&
\end{tikzcd}
\]
（用第III讲1.4。）

j) 设 $\operatorname{Pro}_S\mathscr F$ 是 $\operatorname{Pro}\mathscr F$（第I讲8）中由过渡\Term{morphism}{态射}属于 $S$ 的\Term{pro-object}{亲对象}所确定的\Term{full-subcategory}{满子范畴}。证明\Term{functor}{函子} $Q$ \Term{resolution}{分解}为
\[
\begin{tikzcd}
\mathscr F\arrow[r,hook,"i"]&\operatorname{Pro}_S\mathscr F\arrow[r,"j"]&\mathscr F_{S^{-1}},
\end{tikzcd}
\]
其中 $j$ 是 $\operatorname{Pro}Q$ 在 $\operatorname{Pro}_S\mathscr F$ 上的\Term{restriction}{限制}。证明\Term{functor}{函子} $j$ 有\Term{fully-faithful}{全忠实}\Term{left-adjoint}{左伴随}
\[
A:\mathscr F_{S^{-1}}\longrightarrow\operatorname{Pro}_S\mathscr F.
\]
证明：对 $\mathscr F$ 的每个\Term{object}{对象} $X$，$AQ(X)$ 是\Term{pro-object}{亲对象}
\[
\bigl(Y\longrightarrow X\in S(X)\bigr)\longmapsto Y.
\]
\end{SGAresult}

\begin{SGAresult}{命题}{6.7}
设 $\mathscr E$ 是\Term{cofiltered}{余滤过}\Term{category}{范畴}（第I讲8.7），$\mathscr F\to\mathscr E$ 是\Term{fibered-category}{纤维范畴}（6.1.3）。

1) 设 $\xi$ 是 $\mathscr E$ 的\Term{object}{对象}，$\mathscr F_{/\xi}$ 是经\Term{base-change}{基变换} $\mathscr E_{/\xi}\to\mathscr E$ 得到的 $\mathscr E_{/\xi}$ 上的\Term{fibered-category}{纤维范畴}。\Term{canonical}{典范}\Term{functor}{函子}
\[
\varinjlim_{\mathscr E_{/\xi}}\mathscr F_{/\xi}\longrightarrow\varinjlim_{\mathscr E}\mathscr F
\]
是\Term{category}{范畴}\Term{equivalence}{等价}。
\SourcePage{183}
2) 更一般地，设 $\mathscr E'\to\mathscr E$ 是\Term{cofinal}{余尾}\Term{functor}{函子}（第I讲8）。设 $\mathscr F'\to\mathscr E'$ 是由 $\mathscr F\to\mathscr E$ 经\Term{base-change}{基变换} $\mathscr E'\hookrightarrow\mathscr E$ 得到的\Term{fibered-category}{纤维范畴}。\Term{canonical}{典范}\Term{functor}{函子}
\[
\varinjlim_{\mathscr E'}\mathscr F'\longrightarrow\varinjlim_{\mathscr E}\mathscr F
\]
是\Term{category}{范畴}\Term{equivalence}{等价}。
\end{SGAresult}
\begin{proof}
习题。
\end{proof}

\begin{SGAresult}{习题}{6.8}
1) 设 $g:\mathscr E^{\circ}\to\mathrm{Cat}$ 是\Term{functor}{函子}，$\mathscr G\to\mathscr E$ 是相应的\Term{fibered-category}{纤维范畴}[5]。证明存在\Term{canonical}{典范}\Term{functor}{函子}
\[
\varinjlim_{\mathscr E^{\circ}}\mathscr G\longrightarrow\varinjlim_{\mathscr E^{\circ}}g.
\]
证明此\Term{functor}{函子}不一定是\Term{category}{范畴}\Term{equivalence}{等价}。（可以取 $g$ 为值等于 $\mathrm{Cat}$ 的\Term{final-object}{终对象}的常\Term{functor}{函子}，$\mathscr E$ 为群所关联的\Term{category}{范畴}。）

2) 假定\Term{category}{范畴} $\mathscr E^{\circ}$ \Term{pseudo-filtered}{伪滤过}。证明这时\Term{canonical}{典范}\Term{functor}{函子} $\varinjlim_{\mathscr E^{\circ}}\mathscr G\to\varinjlim_{\mathscr E^{\circ}}g$ 是\Term{category}{范畴}\Term{equivalence}{等价}。
\end{SGAresult}

\begin{SGAresult}{命题}{6.9}
设 $\mathscr F\to\mathscr E$ 是\Term{fibered-category}{纤维范畴}（6.1）。\Term{functor}{函子} $(\mathrm{Cat})\to(\mathrm{Ens})$
\[
C\longmapsto\operatorname{Hom}_{\mathrm{Cart}/\mathscr E}(C\times\mathscr E,\mathscr F)
\]
由\Term{category}{范畴} $\mathscr{H}\!om_{\mathrm{Cart}/\mathscr E}(\mathscr E,\mathscr F)$ 表示。
\end{SGAresult}
\begin{proof}
设 $C$ 是\Term{category}{范畴}，$F:C\times\mathscr E\to\mathscr F$ 是\Term{cartesian}{笛卡尔} $\mathscr E$-\Term{functor}{函子}（6.1.4），$X$ 是 $C$ 的\Term{object}{对象}。\Term{functor}{函子} $\xi\mapsto F((X,\xi))$ 是从 $\mathscr E$ 到 $\mathscr F$ 的\Term{cartesian-functor}{笛卡尔函子}，记作 $F'(X)$，且\Term{functor}{函子}性地依赖于 $X$。于是得到映射
\[
F\longmapsto F':\operatorname{Hom}_{\mathrm{Cart}/\mathscr E}(C\mathscr E,\mathscr F)
\longrightarrow\operatorname{Hom}\bigl(C,\mathscr{H}\!om_{\mathrm{Cart}/\mathscr E}(\mathscr E,\mathscr F)\bigr),
\]
它关于 $C$ 是\Term{functor}{函子}性的，且立即验证它是\Term{bijection}{双射}。
\end{proof}

\textbf{6.10.} \Term{category}{范畴} $\mathscr{H}\!om_{\mathrm{Cart}/}(\mathscr E,\mathscr F)$ 称为\Term{fibered-category}{纤维范畴} $\mathscr F$ 的\Term{cartesian-section}{笛卡尔截面}\Term{category}{范畴}。它有时也称为 $\mathscr F$ 沿 $\mathscr E^{\circ}$ 的\Term{inverse-limit}{射影极限}\Term{category}{范畴}，这时记作 $\varprojlim_{\mathscr E^{\circ}}\mathscr F$ [2]。

\textbf{6.11.} 选择\Term{fibered-category}{纤维范畴} $\mathscr F\to\mathscr E$ 的一个\Term{splitting}{分裂}（scindage），即对 $\mathscr E$ 的每条\Term{arrow}{箭头} $f:\xi\to\eta$，选择\Term{base-change}{基变换}\Term{functor}{函子} $f^*:\mathscr F_\eta\to\mathscr F_\xi$。对 $\mathscr E$ 的每个\Term{object}{对象} $\xi$，以
\[
v_\xi:\varprojlim_{\mathscr E^{\circ}}\mathscr F\longrightarrow\mathscr F_\xi
\]
\SourcePage{184}
表示在 $\xi$ 处的\Term{evaluation-functor}{求值函子}。对每个\Term{morphism}{态射} $f:\xi\to\eta$，有准确到\Term{canonical}{典范}\Term{isomorphism}{同构}交换的\Term{functor}{函子}图：
\[
\begin{tikzcd}
\varprojlim_{\mathscr E^{\circ}}\mathscr F\arrow[r,"v_\xi"]\arrow[dr,"v_\eta"']&\mathscr F_\xi\arrow[d,"f^*"]\\
&\mathscr F_\eta.
\end{tikzcd}
\]
因此\Term{category}{范畴} $\varprojlim_{\mathscr E^{\circ}}\mathscr F$ 表现为 $\xi\mapsto\mathscr F_\xi$ 在\Term{pseudofunctor}{伪函子}意义下的\Term{inverse-limit}{射影极限}。即使 $\xi\mapsto\mathscr F_\xi$ 是真正的\Term{functor}{函子}（即所选的\Term{splitting}{分裂}是\Term{cleavage}{劈裂}（clivage）[5]），\Term{category}{范畴} $\varprojlim_{\mathscr E^{\circ}}\mathscr F$（6.10）与 $\varprojlim_{\mathscr E^{\circ}}\mathscr F_\xi$（第I讲2）一般也不\Term{equivalence}{等价}。
